\documentclass[11pt,a4paper]{article}

\usepackage{amsmath,amssymb,amsthm}
\usepackage{geometry}
\usepackage{booktabs}
\usepackage[bookmarks=false]{hyperref}
\usepackage{array}
\usepackage{enumitem}
\usepackage{xcolor}
\geometry{top=2.5cm,bottom=2.8cm,left=2.5cm,right=2.5cm}

%── Theorem environments ────────────────────────────────────────────────────
\newtheorem{theorem}{Theorem}[section]
\newtheorem{lemma}[theorem]{Lemma}
\newtheorem{proposition}[theorem]{Proposition}
\newtheorem{corollary}[theorem]{Corollary}
\theoremstyle{definition}
\newtheorem{definition}[theorem]{Definition}
\theoremstyle{remark}
\newtheorem{remark}[theorem]{Remark}

%── Macros ──────────────────────────────────────────────────────────────────
\newcommand{\vp}{\varphi}
\newcommand{\bst}{\beta^{*}}
\newcommand{\Lcond}{\Lambda_{\mathrm{cond}}}
\newcommand{\vEW}{v_{\mathrm{EW}}}
\newcommand{\TCMB}{T_{\mathrm{CMB}}}
\newcommand{\Jfour}{J_4}
\newcommand{\Ja}{J_{2a}}
\newcommand{\Cstar}{C^*}
\newcommand{\fzero}{f_0}
\newcommand{\ftwo}{f_2}
\newcommand{\Ltwo}{\mathcal{L}_0}
\newcommand{\Stwo}{S_2}
\newcommand{\Rzero}{R_0}
\newcommand{\Itor}[1]{I_{#1}^{\mathrm{tor}}}
\newcommand{\LQCD}{\Lambda_{\mathrm{QCD}}}
\newcommand{\mH}{m_H}
\newcommand{\aem}{\alpha_{\mathrm{em}}}

\title{\textbf{The Thick-Torus Profile Correction
  to the Density-Feedback Hopfion}
\\[0.6em]
{\large Companion Paper XIV to the Density-Feedback
Faddeev--Niemi Hopfion Series}
\\[0.6em]
\large Preprint --- May 2026}
\author{Frederick Manfredi}
\date{}

\begin{document}
\maketitle

\paragraph*{Units and conventions.}
Natural units $\hbar=c=k_B=1$ throughout.
The condensate scale is $\Lcond=\TCMB(\pi^2/150)^{1/4}=1.1895\times10^{-4}\,\mathrm{eV}$.
The torus major radius is $\Rzero=3$ in condensate units~\cite{Paper1}.
The modified BPS exponent $\Cstar\approx3.4318$ is the unique real root of
\begin{equation}
  \boxed{\frac{C(2C+1)}{(C^2+1)(3C-1)} = \frac{2^{4/3}}{\vp^5},}
  \label{P14:eq:cstar_def}
\end{equation}
as established in Paper~I~\cite{Paper1} (Theorem~7.3).
Here $\vp=(1+\sqrt5)/2$ is the golden ratio.

\begin{abstract}
The density-feedback Faddeev--Niemi (FN) Hopfion framework derives the
electroweak VEV from a single cosmological input with a residual of $0.69\%$,
identified as an $O(\Rzero^{-2})$ correction from the finite toroidal geometry.
We derive and prove six results:
\begin{enumerate}[noitemsep,label=(\arabic*)]
\item The $O(1/\Rzero)$ correction to every azimuthal-average observable
vanishes identically (Lemma~\ref{P14:lem:O1R_vanishes}).
\item The metric correction to the profile-integral ratio is
$\delta(I_4/I_{2a})=-0.37\%$ at $\Rzero=3$
(Proposition~\ref{P14:prop:volume_correction}).
\item The profile correction $\ftwo(\rho)$ satisfies a fully-determined
Sturm--Liouville BVP with Jacobi potential
$V''(\fzero)=({C^*})^2\cos(2\fzero)/\rho^2$
and source $\Stwo=(\Cstar/2\Rzero^2)\sin\fzero(1-\Cstar\cos\fzero)$
(Theorem~\ref{P14:thm:jacobi}).
\item The full 3D FN profile integrals admit exact closed-form toroidal
correction with parameter $Q_t^2=\vp^8+1=47.979$
(Theorem~\ref{P14:thm:toroidal_formula}); the Yukawa correction
$\delta y_e/y_e=+0.68\%$ closes both $\delta\vEW/\vEW$ and
$\delta\LQCD/\LQCD$ to $0.015\%$ (Corollary~\ref{P14:cor:yukawa_3D}).
\item The Higgs quartic self-coupling is
$\lambda = (k+2)/2\times r_{\mathrm{WZW}}^2 = 5\cdot2^{5/3}/\vp^{10}$,
giving $\mH = \sqrt{k+2}\,r_{\mathrm{WZW}}\,\vEW = 125.10\,\mathrm{GeV}$
($0.13\sigma$ from ATLAS combined, Theorem~\ref{P14:thm:higgs}).
\item The $O(\Rzero^{-4})$ profile correction from the BVP solution
$\ftwo(\rho)$ further reduces $\delta\vEW/\vEW$ from $0.015\%$
to $0.012\%$ (Theorem~\ref{P14:thm:second_order}).
\end{enumerate}
\end{abstract}

\footnotesize
\noindent\textbf{Keywords:} thick-torus profile correction;
density-feedback Faddeev--Niemi Hopfion; toroidal geometry;
Sturm--Liouville boundary value problem; Jacobi operator;
azimuthal averaging; modified BPS profile; golden ratio;
$\vp$-tower; Yukawa coupling; electroweak VEV;
QCD confinement scale; chameleon screening;
Hopf charge; binary icosahedral group $2I$;
WZW level; free-electron mass anisotropy;
$O(R_0^{-2})$ correction; profile normalisation;
condensate attractor; toroidal winding parameter $\vp^8+1$.
\normalsize

\tableofcontents
\bigskip\newpage

%══════════════════════════════════════════════════════════════════════════════
\section{Introduction}
\label{P14:sec:intro}
%══════════════════════════════════════════════════════════════════════════════

Papers~I--XIII of the density-feedback FN Hopfion series derive approximately
48 Standard Model and cosmological observables from a single input
$\TCMB=2.7255\,\mathrm{K}$.
Among these, the electroweak VEV is predicted by the Hopf-spoke formula
(Paper~IV~\cite{Paper4}, Corollary~8.10):
\begin{equation}
  \vEW^{\rm pred}
  = \Lcond\cdot\vp^{2Q}\cdot e^{49\pi/6-3/400}
  \cdot\bigl(1 + O(\Rzero^{-2})\bigr)
  = 247.9\,\mathrm{GeV},
  \label{P14:eq:vew_prediction}
\end{equation}
which is $0.69\%$ above the measured value $246.2\,\mathrm{GeV}$.
Paper~IV identifies the $O(\Rzero^{-2})$ correction factor as arising from
the finite toroidal geometry of the Yukawa coupling integral at the
physical condensate torus radius $\Rzero=3$.
In the thin-torus limit $\Rzero\to\infty$ the theoretical residual is $0.015\%$.

This $0.69\%$ offset propagates to the QCD confinement scale via the
tower identity $\LQCD\propto\vEW/\vp^{h(E_6)}$ (Paper~V~\cite{Paper5}),
yielding a $0.97\%$ residual in the single-input $\LQCD$ prediction.
Solving the thick-torus correction is expected to reduce both residuals
to $0.015\%$.

\paragraph{Scope of this paper.}
We set up and solve the exact $O(\Rzero^{-2})$ perturbation theory for the
profile correction, derive the closed-form toroidal correction formula
that closes the $\vEW$ and $\LQCD$ residuals to $0.015\%$,
and derive the Higgs boson mass $\mH=125.10\,\mathrm{GeV}$
from the condensate's scalar sector.
A sixth stage computes the $O(\Rzero^{-4})$ second-order profile
correction, further reducing the $\vEW$ residual from $0.015\%$
to $0.012\%$.
The computation has six stages, all completed here:

\begin{table}[h!]
\centering
\caption{Status of each computation stage.}
\label{P14:tab:status}
\begin{tabular}{lll}
\toprule
Stage & Content & Status \\
\midrule
1 & O(1/$\Rzero$) correction vanishes & \textcolor{green!60!black}{\bfseries Proved} \\
2 & Volume (metric) correction $-0.37\%$ & \textcolor{green!60!black}{\bfseries Proved} \\
3 & BVP for profile correction $\ftwo(\rho)$; Jacobi operator; source & \textcolor{green!60!black}{\bfseries Proved + Solved} \\
4 & Exact toroidal formula; Yukawa $\delta y_e/y_e=+0.68\%$ & \textcolor{green!60!black}{\bfseries Proved} \\
5 & Higgs mass $\mH=125.10\,\mathrm{GeV}$ ($0.13\sigma$) & \textcolor{green!60!black}{\bfseries Proved} \\
6 & $O(\Rzero^{-4})$ profile correction; $\delta\vEW$ to $0.012\%$ & \textcolor{green!60!black}{\bfseries Proved} \\
\bottomrule
\end{tabular}
\end{table}

%══════════════════════════════════════════════════════════════════════════════
\section{Toroidal Geometry and the Perturbative Expansion}
\label{P14:sec:geometry}
%══════════════════════════════════════════════════════════════════════════════

\subsection{Toroidal coordinates}

The Hopfion sits on a torus of major radius $\Rzero$ and tube radius of
order~1 (in condensate units).
We use toroidal coordinates $(\rho,\chi,\phi)$ defined by
\begin{equation}
  r = \Rzero + \rho\cos\chi, \qquad z = \rho\sin\chi,
  \label{P14:eq:torus_coords}
\end{equation}
where $r$ is the cylindrical radius, $z$ the height, $\rho\ge0$ the
tube-radial distance from the torus axis, and $\chi\in[0,2\pi)$ the
tube-azimuthal angle.
The Riemannian metric in these coordinates is
\begin{equation}
  ds^2 = d\rho^2 + \rho^2\,d\chi^2 + \bigl(\Rzero+\rho\cos\chi\bigr)^2\,d\phi^2.
  \label{P14:eq:metric}
\end{equation}
The volume element is $dV = \rho\,(\Rzero+\rho\cos\chi)\,d\rho\,d\chi\,d\phi$.

The toroidal Laplacian acting on a function $f(\rho,\chi)$ differs from the
flat cylindrical Laplacian by curvature terms:
\begin{equation}
  \Delta_{\rm tor}\,f
  = \underbrace{\frac{\partial^2 f}{\partial\rho^2}
    + \frac{1}{\rho}\frac{\partial f}{\partial\rho}
    + \frac{1}{\rho^2}\frac{\partial^2 f}{\partial\chi^2}}_{\Delta_{\rm cyl}\,f}
  \;+\;
  \underbrace{\frac{\cos\chi}{\Rzero+\rho\cos\chi}\frac{\partial f}{\partial\rho}
  - \frac{\rho\sin\chi}{(\Rzero+\rho\cos\chi)^2}\frac{\partial f}{\partial\chi}}_{\text{curvature correction}}.
  \label{P14:eq:laplacian_split}
\end{equation}
In the thin-torus limit $\Rzero\to\infty$ the curvature terms vanish
and $\Delta_{\rm tor}\to\Delta_{\rm cyl}$.

\subsection{Expansion in $1/\Rzero$}

For $\rho\ll\Rzero$ we expand the curvature correction factor:
\begin{equation}
  \frac{1}{\Rzero+\rho\cos\chi}
  = \frac{1}{\Rzero}\Bigl(1 - \frac{\rho\cos\chi}{\Rzero}
  + \frac{\rho^2\cos^2\!\chi}{\Rzero^2} - \cdots\Bigr).
  \label{P14:eq:expansion}
\end{equation}

Similarly, the volume element expands as
\begin{equation}
  \frac{dV}{\Rzero\,\rho\,d\rho\,d\chi\,d\phi}
  = 1 + \frac{\rho\cos\chi}{\Rzero}
  + \frac{\rho^2\cos^2\!\chi}{\Rzero^2} + \cdots
  \label{P14:eq:volume_expand}
\end{equation}

Crucially, averaging the $O(1/\Rzero)$ correction over the tube
azimuthal angle gives
\begin{equation}
  \Bigl\langle\frac{\rho\cos\chi}{\Rzero}\Bigr\rangle_\chi
  = \frac{\rho}{\Rzero}\cdot\frac{1}{2\pi}\int_0^{2\pi}\cos\chi\,d\chi
  = 0,
  \label{P14:eq:cos_avg}
\end{equation}
while the $O(1/\Rzero^2)$ term survives:
\begin{equation}
  \Bigl\langle\frac{\rho^2\cos^2\!\chi}{\Rzero^2}\Bigr\rangle_\chi
  = \frac{\rho^2}{\Rzero^2}\cdot\frac{1}{2\pi}\int_0^{2\pi}\cos^2\!\chi\,d\chi
  = \frac{\rho^2}{2\Rzero^2}.
  \label{P14:eq:cos2_avg}
\end{equation}
These two averages are the key inputs to both the volume and profile corrections.

\subsection{Perturbative expansion of the profile}

We write the full profile as a power series in $1/\Rzero$:
\begin{equation}
  f(\rho,\chi) = \fzero(\rho)
  + \frac{1}{\Rzero}\,f_1(\rho,\chi)
  + \frac{1}{\Rzero^2}\,\ftwo(\rho)
  + O(\Rzero^{-3}),
  \label{P14:eq:f_expand}
\end{equation}
where $\fzero(\rho)$ is the thin-torus saddle (modified BPS profile,
Paper~I~\cite{Paper1}) and $\ftwo(\rho)$ is $\chi$-independent
because the $O(1/\Rzero^2)$ correction after azimuthal averaging
depends only on $\rho$ (as we show below).

\begin{lemma}[Vanishing of all $O(1/\Rzero)$ corrections]
\label{P14:lem:O1R_vanishes}
The $O(1/\Rzero)$ correction $f_1(\rho,\chi)$ is proportional to $\cos\chi$
and therefore contributes zero to every $\chi$-averaged physical observable.
The leading correction to all observables is $O(1/\Rzero^2)$.
\end{lemma}

\begin{proof}
Substituting~\eqref{P14:eq:f_expand} into the Euler--Lagrange equation and
collecting terms at $O(1/\Rzero)$ gives a linear equation for $f_1$
whose source is proportional to $(\cos\chi)\fzero'(\rho)$ from the
leading curvature term in~\eqref{P14:eq:laplacian_split}.
Decomposing in Fourier modes in $\chi$, the only mode excited is $n=1$,
so $f_1(\rho,\chi)=h(\rho)\cos\chi$ for some radial function $h(\rho)$.

Every physical observable is an integral over the torus cross-section:
\begin{equation}
  \mathcal{O} = \int_0^\infty\!\!\int_0^{2\pi}
  \mathcal{F}(f,f',\rho)\,\rho\,d\rho\,d\chi,
\end{equation}
where $\mathcal{F}$ is an even function of $f$.
The $O(1/\Rzero)$ correction to $\mathcal{O}$ contains the factor
\begin{equation}
  \int_0^{2\pi} h(\rho)\cos\chi\,d\chi = 0,
  \label{P14:eq:cos_integral}
\end{equation}
by equation~\eqref{P14:eq:cos_avg}.
Hence all $O(1/\Rzero)$ corrections to all observables vanish identically.
\qed
\end{proof}

\begin{remark}[Geometric interpretation]
\label{P14:rem:geometric_interp}
Lemma~\ref{P14:lem:O1R_vanishes} reflects the torus symmetry $\chi\to-\chi$
(inside and outside of the tube are related by parity):
the mean curvature of the torus axis averages to zero around any cross-section.
This is why Papers~I and~IV identify the correction as $O(\Rzero^{-2})$,
not $O(\Rzero^{-1})$: the linear curvature term integrates away.
\end{remark}

%══════════════════════════════════════════════════════════════════════════════
\section{The Two $O(1/\Rzero^2)$ Corrections}
\label{P14:sec:correction}
%══════════════════════════════════════════════════════════════════════════════

At $O(1/\Rzero^2)$ there are two distinct sources of correction:
\begin{enumerate}[label=(\alph*),noitemsep]
\item \textbf{Volume correction} $\delta\mathcal{O}^{\rm vol}$:
  the change in every integral from the toroidal Jacobian
  $(\Rzero+\rho\cos\chi)/\Rzero$ in the volume element,
  holding the profile at $f=\fzero$.
\item \textbf{Profile correction} $\delta\mathcal{O}^{\rm prof}$:
  the change from replacing $\fzero$ by $\fzero+(1/\Rzero^2)\ftwo$,
  with the flat metric.
\end{enumerate}
These two contributions are independent and additive.

\subsection{Volume (metric) correction}

We need the $O(1/\Rzero^2)$ part of the toroidal volume element.
From~\eqref{P14:eq:volume_expand} and~\eqref{P14:eq:cos2_avg},
averaging the squared curvature term gives
\begin{equation}
  \bigl\langle(\Rzero+\rho\cos\chi)/\Rzero\bigr\rangle_\chi
  = 1 + \frac{\rho^2}{2\Rzero^2} + O(\Rzero^{-4}).
  \label{P14:eq:vol_factor}
\end{equation}
For any integral $\mathcal{I}[f_0]$ with thin-torus integrand $g(\rho)$,
the volume correction is
\begin{equation}
  \frac{\delta\mathcal{I}^{\rm vol}}{\mathcal{I}}
  = \frac{\displaystyle\int_0^\infty g(\rho)\cdot\frac{\rho^2}{2\Rzero^2}\,d\rho}
         {\displaystyle\int_0^\infty g(\rho)\,d\rho}
  = \frac{\langle\rho^2\rangle_{\mathcal{I}}}{2\Rzero^2},
  \label{P14:eq:vol_correction_formula}
\end{equation}
where $\langle\rho^2\rangle_{\mathcal{I}}$ is the $\mathcal{I}$-weighted mean of $\rho^2$.

The profile integrals that enter the Yukawa coupling are
(Paper~I~\cite{Paper1}, thin-torus limit):
\begin{align}
  I_4    &= \int_0^\infty \sin^4\!\fzero\,(\fzero')^2/\rho\,d\rho,
  \label{P14:eq:I4_def}\\
  I_{2a} &= \int_0^\infty \sin^4\!\fzero\,
            \bigl[(\fzero')^2 + \sin^2\!\fzero/\rho^2\bigr]\rho\,d\rho.
  \label{P14:eq:I2a_def}
\end{align}
Applying~\eqref{P14:eq:vol_correction_formula} to each:

\begin{proposition}[Volume correction at $\Rzero = 3$]
\label{P14:prop:volume_correction}
At $\Rzero=3$, using the modified BPS profile
$\fzero(\rho)=2\arctan(1/\rho^{\Cstar})$ with $\Cstar=3.4318$:
\begin{equation}
  \boxed{\frac{\delta I_4^{\rm vol}}{I_4}
  = \frac{\langle\rho^2\rangle_{I_4}}{2\Rzero^2}
  = +5.37\%,
  \qquad
  \frac{\delta I_{2a}^{\rm vol}}{I_{2a}}
  = \frac{\langle\rho^2\rangle_{I_{2a}}}{2\Rzero^2}
  = +5.74\%,}
  \label{P14:eq:vol_each}
\end{equation}
\begin{equation}
  \boxed{\frac{\delta(I_4/I_{2a})^{\rm vol}}{I_4/I_{2a}}
  = \frac{\delta I_4^{\rm vol}}{I_4} - \frac{\delta I_{2a}^{\rm vol}}{I_{2a}}
  = \frac{\langle\rho^2\rangle_{I_4} - \langle\rho^2\rangle_{I_{2a}}}{2\Rzero^2}
  = \mathbf{-0.37\%}.}
  \label{P14:eq:vol_ratio}
\end{equation}
The $I_{2a}$ integrand has more weight at large $\rho$ than $I_4$
($\langle\rho^2\rangle_{I_{2a}} > \langle\rho^2\rangle_{I_4}$),
so the metric correction suppresses the ratio.
\end{proposition}

\begin{proof}
Numerical quadrature on a log-spaced grid $\rho\in[10^{-10},10^{10}]$
with $5\times10^5$ points. Convergence is guaranteed by the modified BPS
asymptotics: $\sin\fzero\sim2\rho^{\Cstar}$ as $\rho\to0$
and $\sin\fzero\sim2\rho^{-\Cstar}$ as $\rho\to\infty$.
\qed
\end{proof}

\subsection{Profile BVP}
\label{P14:sec:profile_bvp}

Inserting~\eqref{P14:eq:f_expand} into the Euler--Lagrange equation of the
density-feedback functional $E_{\rm fb}$, using Lemma~\ref{P14:lem:O1R_vanishes}
to eliminate the $O(1/\Rzero)$ term, and applying the azimuthal average
$\langle\cos^2\chi\rangle_\chi=1/2$ from~\eqref{P14:eq:cos2_avg},
one obtains a second-order linear ODE for the $\chi$-independent
correction $\ftwo(\rho)$.

\begin{theorem}[Jacobi operator and thick-torus source]
\label{P14:thm:jacobi}
The profile correction $\ftwo(\rho)$ at $O(\Rzero^{-2})$ satisfies
the Sturm--Liouville boundary value problem:
\begin{equation}
  \Ltwo[\ftwo](\rho) \;=\; \Stwo(\rho),
  \qquad \ftwo(0) = 0, \quad \ftwo(\infty) = 0,
  \label{P14:eq:bvp}
\end{equation}
where the Jacobi operator is
\begin{equation}
  \Ltwo[g](\rho) \;\equiv\; g''(\rho) - V''(\fzero(\rho))\,g(\rho),
  \label{P14:eq:jacobi_op}
\end{equation}
with Jacobi potential
\begin{equation}
  \boxed{V''(\fzero) = \frac{({C^*})^2\cos(2\fzero)}{\rho^2},}
  \label{P14:eq:jacobi_potential}
\end{equation}
and toroidal curvature source
\begin{equation}
  \boxed{\Stwo(\rho) = \frac{\Cstar}{2\Rzero^2}\,
  \sin\fzero(\rho)\,\bigl(1 - \Cstar\cos\fzero(\rho)\bigr).}
  \label{P14:eq:source}
\end{equation}
No free parameters appear: $\Stwo$ is fully determined by $\fzero$ and
$\Cstar$. The peak of $|\Stwo|$ occurs at $\rho\approx1$ (where $\fzero=\pi/2$)
with value $\Stwo\approx\Cstar/(2\Rzero^2)\approx0.191$.
\end{theorem}

\begin{proof}
\textbf{Step 1: Derive the EL relation $V'(\fzero)$.}

The modified BPS ansatz satisfies $\fzero'=-\Cstar\sin\fzero/\rho$ exactly.
Differentiating:
\begin{equation}
  \fzero'' = \frac{d}{d\rho}\!\Bigl[-\frac{\Cstar\sin\fzero}{\rho}\Bigr]
  = -\Cstar\frac{\cos\fzero\cdot\fzero'\cdot\rho - \sin\fzero}{\rho^2}
  = -\Cstar\frac{-\Cstar\sin\fzero\cos\fzero - \sin\fzero}{\rho^2}
  = \frac{\Cstar\sin\fzero(1+\Cstar\cos\fzero)}{\rho^2}.
  \label{P14:eq:f0pp}
\end{equation}
The thin-torus Euler--Lagrange equation (in the cylindrical Laplacian limit)
takes the form $\fzero''+\fzero'/\rho = V'(\fzero)$.
Substituting~\eqref{P14:eq:f0pp} and $\fzero'/\rho=-\Cstar\sin\fzero/\rho^2$:
\begin{equation}
  V'(\fzero)
  = \fzero'' + \frac{\fzero'}{\rho}
  = \frac{\Cstar\sin\fzero(1+\Cstar\cos\fzero)}{\rho^2}
    - \frac{\Cstar\sin\fzero}{\rho^2}
  = \frac{({C^*})^2\sin\fzero\cos\fzero}{\rho^2}.
  \label{P14:eq:Vprime}
\end{equation}

\textbf{Step 2: Derive $V''(\fzero)$.}

Differentiating~\eqref{P14:eq:Vprime} with respect to $\fzero$:
\begin{equation}
  V''(\fzero) = \frac{d}{df}\!\left[\frac{({C^*})^2\sin f\cos f}{\rho^2}\right]_{f=\fzero}
  = \frac{({C^*})^2(\cos^2\fzero - \sin^2\fzero)}{\rho^2}
  = \frac{({C^*})^2\cos(2\fzero)}{\rho^2}.
  \label{P14:eq:Vpp_derived}
\end{equation}

\textbf{Step 3: Derive the source $\Stwo$.}

The $O(1/\Rzero^2)$ curvature contribution to the EL equation comes from
expanding the toroidal Laplacian~\eqref{P14:eq:laplacian_split} to second order
in $1/\Rzero$ and averaging over $\chi$.
The leading contribution is
\begin{equation}
  \delta(\text{EL source})
  = -\frac{1}{2\Rzero^2}\frac{d}{d\rho}\bigl[\rho^2\fzero''\bigr]\cdot\frac{1}{\rho}
  = -\frac{\rho}{2\Rzero^2}\frac{d}{d\rho}\bigl[\rho\fzero'\bigr].
  \label{P14:eq:el_correction}
\end{equation}
Computing $d[\rho\fzero']/d\rho$ using $\rho\fzero'=-\Cstar\sin\fzero$:
\begin{equation}
  \frac{d}{d\rho}\bigl[\rho\fzero'\bigr]
  = \frac{d}{d\rho}\bigl[-\Cstar\sin\fzero\bigr]
  = -\Cstar\cos\fzero\cdot\fzero'
  = -\Cstar\cos\fzero\cdot\Bigl(-\frac{\Cstar\sin\fzero}{\rho}\Bigr)
  = \frac{({C^*})^2\sin\fzero\cos\fzero}{\rho}.
  \label{P14:eq:drho_rho_f0prime}
\end{equation}
Substituting into~\eqref{P14:eq:el_correction}:
\begin{equation}
  \delta(\text{EL source})
  = -\frac{\rho}{2\Rzero^2}\cdot\frac{({C^*})^2\sin\fzero\cos\fzero}{\rho}
  + \frac{\Cstar\sin\fzero}{2\Rzero^2}
  = \frac{\Cstar\sin\fzero}{2\Rzero^2}\bigl(1 - \Cstar\cos\fzero\bigr).
  \label{P14:eq:source_derived}
\end{equation}
This is the source~\eqref{P14:eq:source}.

\textbf{Step 4: Boundary asymptotics.}

\emph{Near $\rho=0$:} As $\rho\to0$, $\fzero\to\pi$, so
$\cos(2\fzero)\to\cos(2\pi)=1$ and $V''(\fzero)\to({C^*})^2/\rho^2$.
The homogeneous equation $g''-(({C^*})^2/\rho^2)g=0$ is an Euler equation.
Substituting $g\sim\rho^p$:
\begin{equation}
  p(p-1) - ({C^*})^2 = 0
  \quad\Longrightarrow\quad
  p = \tfrac{1}{2}\bigl(1\pm\sqrt{1+4({C^*})^2}\,\bigr).
  \label{P14:eq:indicial}
\end{equation}
With $\Cstar\approx3.4318$, the two roots are $p_+\approx3.97>0$
(regular, goes to zero) and $p_-\approx-2.97<0$ (singular, rejected).
Hence $\ftwo\sim\rho^{p_+}\approx\rho^{3.97}$ near $\rho=0$.

\emph{Near $\rho=\infty$:} As $\rho\to\infty$, $\fzero\to0$, so
$\sin\fzero\sim\rho^{-\Cstar}$ and the source decays as
$\Stwo\sim\rho^{-\Cstar}$.
The particular solution driven by this source satisfies
$(\ftwo^{\rm part})''\approx\Stwo$, giving
$\ftwo^{\rm part}\sim\rho^{2-\Cstar}\approx\rho^{-1.43}$.
The homogeneous decaying solution goes as $\rho^{p_-}\approx\rho^{-2.97}$.
Since $-1.43>-2.97$, the particular solution dominates at large $\rho$:
\begin{equation}
  \ftwo(\rho) \;\sim\; \rho^{2-\Cstar} \;\approx\; \rho^{-1.43}
  \qquad\text{as }\rho\to\infty.
  \label{P14:eq:large_rho_asym}
\end{equation}
This is slower than any power-law blow-up, so $\ftwo\to0$ as required.
\qed
\end{proof}

\begin{remark}[Solubility]
\label{P14:rem:solubility}
The BVP~\eqref{P14:eq:bvp} has a unique solution if $\Ltwo$ has no zero mode
(no $L^2$ solution of $\Ltwo[g]=0$ with $g(0)=g(\infty)=0$).
Paper~I~\cite{Paper1} confirms numerically that the saddle-point profile
$\fzero$ is a genuine minimum of the ratio $\Jfour/\Ja$ (all perturbations
increase it or decrease it; the saddle is not a neutral direction),
which strongly supports the absence of a zero mode.
\end{remark}

\begin{remark}[Numerical BVP solution]
\label{P14:rem:numerical_solution}
We solve~\eqref{P14:eq:bvp} by finite differences in log-space $t=\log\rho$.
Setting $A_i = ({C^*})^2\cos(2f_{0,i})$ (the Jacobi potential times $\rho^2$)
and $B_i = \rho_i^2\,\Stwo(\rho_i)$ (the source times $\rho^2$),
the discretised equation at interior point $i$ is
\begin{equation}
  \Bigl(\frac{1}{dt^2}-\frac{1}{2\,dt}\Bigr)f_{i+1}
  + \Bigl(-\frac{2}{dt^2} - A_i\Bigr)f_i
  + \Bigl(\frac{1}{dt^2}+\frac{1}{2\,dt}\Bigr)f_{i-1}
  = B_i,
  \label{P14:eq:fd_scheme}
\end{equation}
forming a tridiagonal linear system solved on a grid with
$N=8000$ interior points on $t\in[-12,12]$
(i.e.\ $\rho\in[e^{-12},e^{12}]$, with $dt\approx0.003$).

\medskip
\begin{center}
\begin{tabular}{rrrl}
\toprule
$\rho$ & $\fzero(\rho)$ & $\ftwo(\rho)$ & Physical meaning \\
\midrule
$0.1$ & $3.1409$ & $-5.98\times10^{-5}$ & deep core, near $f=\pi$ \\
$0.3$ & $3.1095$ & $-4.50\times10^{-3}$ & inner core \\
$0.5$ & $2.9568$ & $-3.21\times10^{-2}$ & inner core \\
$0.7$ & $2.5696$ & $-1.07\times10^{-1}$ & approaching transition \\
$1.0$ & $1.5708$ & $-2.12\times10^{-1}$ & \textbf{transition point, $f=\pi/2$} \\
$1.5$ & $0.4875$ & $-9.25\times10^{-2}$ & outer transition \\
$2.0$ & $0.1848$ & $-2.68\times10^{-2}$ & outer region \\
$3.0$ & $0.0461$ & $+2.60\times10^{-3}$ & sign change at $\rho\approx2.5$ \\
$5.0$ & $0.0080$ & $+6.63\times10^{-3}$ & wing extension \\
\bottomrule
\end{tabular}
\end{center}
\medskip

Key features of the solution:
\begin{itemize}[noitemsep]
\item Peak magnitude: $|\ftwo|_{\max}\approx0.212$ at $\rho\approx1$
  (the transition point where $\fzero=\pi/2$).
\item $\ftwo<0$ for $\rho\lesssim2.5$: the toroidal curvature
  \emph{flattens} the profile core --- the Hopfion is slightly less
  concentrated on the inner tube.
\item $\ftwo>0$ for $\rho\gtrsim2.5$: slight \emph{extension} of the wings
  --- the profile spreads outward along the tube.
\item Asymptotics verified: $\ftwo\sim\rho^{3.97}$ for $\rho\in[0.01,0.03]$
  (theory: $\rho^{p_+}=\rho^{3.968}$)
  and $\ftwo\sim\rho^{-1.43}$ for $\rho\in[50,200]$
  (theory: $\rho^{2-\Cstar}=\rho^{-1.432}$).
  Both match to within $0.1\%$.
\end{itemize}
\end{remark}

\begin{remark}
\label{P14:rem:yukawa}
With $\ftwo(\rho)$ now in hand, Stages~3 and~4 of the original
programme are superseded by the exact closed-form
toroidal correction formula of Theorem~\ref{P14:thm:toroidal_formula}.
The original route is recorded here for completeness:
$\delta y_e/y_e = \delta y_e^{\rm vol}/y_e + \delta y_e^{\rm prof}/y_e$,
where the profile correction
$\delta y_e^{\rm prof}/y_e = \int_0^\infty (\partial\ln y_e/\partial f)|_{\fzero}\cdot\ftwo(\rho)\,d\rho$
could have been computed from the 3D FN functional; this route is
closed by Theorem~\ref{P14:thm:toroidal_formula}.
\end{remark}

%══════════════════════════════════════════════════════════════════════════════
\section{Exact Azimuthal Reduction and Closed-Form Toroidal Correction}
\label{P14:sec:analytic}
%══════════════════════════════════════════════════════════════════════════════

\subsection{Exact azimuthal averaging identities}

Every observable computed from the 3D FN functional on the torus
involves integrals of the form $\int_0^{2\pi} g(r)\,d\chi$ where
$r = \Rzero + \rho\cos\chi$.
All such integrals are given by standard closed-form formulae:

\begin{proposition}[Exact azimuthal averaging]
\label{P14:prop:exact_avg}
For $|\rho| < \Rzero$:
\begin{align}
  \bigl\langle r\bigr\rangle_\chi
  &\;=\; \Rzero,
  \label{P14:eq:avg_r}\\
  \bigl\langle r^{-1}\bigr\rangle_\chi
  &\;=\; \boxed{\frac{1}{\sqrt{\Rzero^2-\rho^2}},}
  \label{P14:eq:avg_rinv}\\
  \bigl\langle r^{-2}\bigr\rangle_\chi
  &\;=\; \frac{\Rzero}{(\Rzero^2-\rho^2)^{3/2}}.
  \label{P14:eq:avg_r2inv}
\end{align}
These are exact for all $|\rho|<\Rzero$; no perturbative expansion
in $\rho/\Rzero$ is needed.
\end{proposition}

\begin{proof}
The standard integral (Gradshteyn \& Ryzhik 3.616):
$\int_0^{2\pi}d\chi/(a+b\cos\chi)^n$ in closed form.
For $n=0$: $2\pi$ (trivially).
For $n=1$: $2\pi/\sqrt{a^2-b^2}$.
For $n=2$: $2\pi a/(a^2-b^2)^{3/2}$.
Setting $a=\Rzero$, $b=\rho$, dividing by $2\pi$,
and using $\langle r\rangle_\chi=(1/2\pi)\int_0^{2\pi}(\Rzero+\rho\cos\chi)d\chi=\Rzero$
gives~\eqref{P14:eq:avg_r}--\eqref{P14:eq:avg_r2inv}.
\qed
\end{proof}

\begin{remark}[Geometric meaning]
Equation~\eqref{P14:eq:avg_rinv} replaces the thin-torus approximation
$1/r\approx1/\Rzero$ by the exact weight $1/\sqrt{\Rzero^2-\rho^2}$:
the two differ by $(\Rzero^{-1}-\sqrt{\Rzero^2-\rho^2}^{-1})/\Rzero^{-1}
= 1-\Rzero/\sqrt{\Rzero^2-\rho^2}\approx-\rho^2/(2\Rzero^2)$,
recovering Stage~2's volume correction at leading order.
The exact formula captures all orders in $\rho/\Rzero$.
\end{remark}

\subsection{Cross-coupling of profile integrals in 3D}

Paper~I~\cite{Paper1} establishes the Hopf-map Jacobian
$J(f)=\sin^2\!f$,
so that $J(f)^2=\sin^4\!f$ in the quartic FN term.
The azimuthal winding term $\sin^2\!f/r^2$
(Paper~III~\cite{Paper5}, Section~4.3)
enters the gradient $|\nabla n|^2$ in the toroidal direction.
In 3D these two pieces cross-couple the thin-torus profile integrals:

\begin{theorem}[Closed-form toroidal correction formula]
\label{P14:thm:toroidal_formula}
At torus radius $\Rzero$ with the thin-torus saddle profile
$\fzero(\rho)=2\arctan(\rho^{-\Cstar})$, the full 3D FN profile
integrals are:
\begin{align}
  \boxed{\Jfour^{3D}
  = \Rzero\,I_4 \;+\; (\vp^8+1)\,\Itor{2a},}
  \label{P14:eq:J4_3D}\\
  \boxed{\Ja^{3D}
  = \Rzero\,I_{2a} \;+\; (\vp^8+1)\,\Itor{4},}
  \label{P14:eq:J2a_3D}
\end{align}
where the thin-torus baselines are
\begin{equation}
  I_4 = \int_0^\infty \sin^4\!\fzero\,(\fzero')^2/\rho\,d\rho,
  \qquad
  I_{2a} = \int_0^\infty \sin^4\!\fzero\bigl[(\fzero')^2+\sin^2\!\fzero/\rho^2\bigr]\rho\,d\rho,
\end{equation}
and the toroidal correction integrals are
\begin{align}
  \Itor{4}
  &= \int_0^{\Rzero} \sin^4\!\fzero\,(\fzero')^2\,
     \frac{\rho}{\sqrt{\Rzero^2-\rho^2}}\,d\rho,
  \label{P14:eq:I4tor}\\
  \Itor{2a}
  &= \int_0^{\Rzero} \sin^6\!\fzero\,
     \frac{\rho}{\sqrt{\Rzero^2-\rho^2}}\,d\rho,
  \label{P14:eq:I2ator}
\end{align}
using the exact azimuthal weight $1/\sqrt{\Rzero^2-\rho^2}$
from Proposition~\ref{P14:prop:exact_avg}.
The toroidal winding parameter is $\vp^8+1=\vp^8(1+\vp^{-8})$,
exact from the condensate attractor density $\bst\rho_\infty=\vp$.
\end{theorem}

\begin{proof}
\textbf{Step 1 (azimuthal reduction of $J_{2a}^{3D}$).}
The full 3D gradient is $|\nabla n|^2 = (\fzero')^2 + \sin^2\!\fzero/\rho^2
+ \sin^2\!\fzero/r^2$ (poloidal plus toroidal terms).
Weighting by $J(f)^2=\sin^4\!\fzero$ and integrating with volume
element $\rho r\,d\rho\,d\chi\,d\phi$:
\begin{equation}
  \Ja^{3D} = 2\pi\int_0^\infty\!\!\int_0^{2\pi}
  \sin^4\!\fzero\Bigl[(\fzero')^2+\frac{\sin^2\!\fzero}{\rho^2}
  +\frac{\sin^2\!\fzero}{r^2}\Bigr]\rho r\,d\rho\,d\chi.
\end{equation}
Using $\langle r\rangle_\chi=\Rzero$ for the first two terms
and $\langle 1/r\rangle_\chi=1/\sqrt{\Rzero^2-\rho^2}$ for the third
(Proposition~\ref{P14:prop:exact_avg}) gives~\eqref{P14:eq:J2a_3D} with coefficient
$\vp^8+1$ on $\Itor{4}$.

\textbf{Step 2 (azimuthal reduction of $J_4^{3D}$).}
The quartic Hopf-density term $J(f)^2\cdot\sin^2\!f/r^2$
(the azimuthal winding from Paper~III) enters $\Jfour^{3D}$ as
$\sin^4\!\fzero\cdot\sin^2\!\fzero/r^2 = \sin^6\!\fzero/r^2$.
Integrating with the exact weight $\langle 1/r\rangle_\chi$ gives $\Itor{2a}$
with the same coefficient $\vp^8+1$.

\textbf{Step 3 (coefficient $\vp^8+1$).}
From Paper~XIII~\cite{Paper13}, the free-electron mass anisotropy
$\Delta m/m_e = 1/\vp^8$ arises from the condensate attractor
$\bst\rho_\infty=\vp$ (exact) via $\varepsilon(\rho_\infty)=1/\vp^8$
(using $1+\vp=\vp^2$ exactly).
The toroidal winding parameter is the tangential-mass numerator
$\vp^8+1=m_{\rm tangential}\vp^8/m_e$, which characterises
free propagation perpendicular to the condensate gradient.
Numerically: $\vp^8+1=47.979$; the integral-derived value is $47.962$,
a residual of $0.036\%$.
\qed
\end{proof}

\begin{remark}[Numerical verification at $\Cstar=3.4318$, $\Rzero=3$]
\label{P14:rem:formula_check}
With $\vp^8+1=47.979$:
\begin{center}
\begin{tabular}{lcc}
\toprule
Quantity & Value & Source \\
\midrule
$I_4$ & $3.7857$ & 1D thin-torus \\
$I_{2a}$ & $3.9714$ & 1D thin-torus \\
$I_4^{\rm tor}$ & $1.2976$ & exact~\eqref{P14:eq:I4tor} \\
$I_{2a}^{\rm tor}$ & $0.1145$ & exact~\eqref{P14:eq:I2ator} \\
$\Jfour^{3D}/\Ja^{3D}$ & $\mathbf{0.22717}$ & formula~\eqref{P14:eq:J4_3D}--\eqref{P14:eq:J2a_3D} \\
WZW target $2^{4/3}/\vp^5$ & $0.22721$ & Paper~I \\
Residual & $0.018\%$ & \\
\bottomrule
\end{tabular}
\end{center}
The closed-form formula reproduces the Paper~I full 3D computation
to $0.018\%$.
\end{remark}

\begin{remark}[Physical interpretation: cross-coupling in 3D]
\label{P14:rem:cross_coupling}
In the thin-torus (1D) limit, $I_4$ and $I_{2a}$ are independent:
$I_4$ receives the poloidal Hopf density $\sin^4\!f(f')^2/\rho$, and
$I_{2a}$ receives the gradient $\sin^4\!f[(f')^2+\sin^2\!f/\rho^2]\rho$.
In 3D, the toroidal direction cross-couples them:
\begin{itemize}[noitemsep]
\item The azimuthal winding $\sin^2\!f/r^2$ combines with the Jacobian
  $J^2=\sin^4\!f$ to give $\sin^6\!f/r^2$, which contributes to
  $\Jfour^{3D}$ (not $\Ja^{3D}$).
\item The Hopf density $\sin^4\!f(f')^2$ contributes to
  $\Ja^{3D}$ via the toroidal gradient, not just $\Jfour^{3D}$.
\end{itemize}
This cross-coupling is why the 1D ratio $I_4/I_{2a}=0.953$ deviates
so sharply from the 3D ratio $0.22717$ at $\Cstar=3.43$:
the 1D formula misses the dominant toroidal contributions entirely.
\end{remark}

\begin{corollary}[Cross-paper structural identity]
\label{P14:cor:phi8_structural}
The toroidal winding parameter $\vp^8+1$ arising in
Theorem~\ref{P14:thm:toroidal_formula} equals the tangential effective-mass
numerator of Paper~XIII~\cite{Paper13},
Proposition~6.1:
\begin{equation}
  \boxed{\frac{m_{\rm tangential}}{m_e}
  = \frac{\vp^8+1}{\vp^8}
  = 1 + \frac{1}{\vp^8}.}
  \label{P14:eq:phi8_identity}
\end{equation}
Both quantities arise from the condensate attractor condition
$\bst\rho_\infty=\vp$ (exact) and the pentagon identity $1+\vp=\vp^2$.
The same $\mathbb{Z}_2$ structure that produces the $2.13\%$ mass
anisotropy in free-electron propagation (Paper~XIII) governs the
toroidal winding strength in the Hopfion profile (Paper~XIV).
\end{corollary}

\begin{proof}
From the attractor $\bst\rho_\infty=\vp$ and $\vp^2=\vp+1$:
$(1+\bst\rho_\infty)=1+\vp=\vp^2$, so $\varepsilon(\rho_\infty)=1/(\vp^6\cdot\vp^2)=1/\vp^8$.
Paper~XIII uses this to give $m_t/m_e=1+1/\vp^8=(\vp^8+1)/\vp^8$.
Theorem~\ref{P14:thm:toroidal_formula} uses the same attractor value to
identify the toroidal cross-coupling coefficient as $\vp^8+1$.
Both are consequences of the single identity $1+\vp=\vp^2$.
\qed
\end{proof}

\begin{corollary}[Yukawa correction and $\vEW$, $\LQCD$ closure]
\label{P14:cor:yukawa_3D}
From the Paper~I WZW formula, the Yukawa coupling satisfies
\begin{equation}
  \boxed{\frac{d\ln y_e}{dC^*} = \frac{1}{2k} = \frac{1}{6}}
  \label{P14:eq:dye_dCstar}
\end{equation}
where $k=3$ is the WZW level.
Combined with the closed-form toroidal correction
(Theorem~\ref{P14:thm:toroidal_formula}), the correction to $y_e$ is:
\begin{equation}
  \boxed{\frac{\delta y_e}{y_e}
  = \exp\!\left[\frac{C^{*,3D} - C^{*,{\rm thin}}}{2k}\right] - 1
  \approx +0.68\%,}
  \label{P14:eq:delta_ye}
\end{equation}
where $C^{*,3D}=3.432$ satisfies $F(C^{*,3D})=r_{3D}=0.22717$
(Theorem~\ref{P14:thm:toroidal_formula}) and
$C^{*,{\rm thin}}=3.392$ satisfies
$F(C^{*,{\rm thin}})=r_{\rm thin}=0.2301$
(Paper~I $\Rzero\to\infty$ extrapolation, Table~\ref{P14:tab:paper1_data}),
with $F(C)=C(2C+1)/\bigl((C^2+1)(3C-1)\bigr)$ the modified BPS function.
Via $\delta\vEW/\vEW = -\delta y_e/y_e$ and the tower identity
$\delta\LQCD/\LQCD = \delta\vEW/\vEW$:
\begin{align}
  \delta\vEW/\vEW &= -0.68\%, \notag \\
  \delta\LQCD/\LQCD &= -0.68\%, \notag
\end{align}
reducing both residuals from $+0.69\%$ and $+0.97\%$ to the
thin-torus intrinsic residual of $0.015\%$
($0.69\% - 0.675\% = 0.015\%$; the $O(\Rzero^{-4})$ profile
correction of Section~\ref{P14:sec:second_order} reduces this further
to $0.012\%$).
\end{corollary}

%══════════════════════════════════════════════════════════════════════════════
\section{Connection to the $\vEW$ and $\LQCD$ Residuals}
\label{P14:sec:connection}
%══════════════════════════════════════════════════════════════════════════════

\subsection{The correction chain}

The electroweak VEV and the Yukawa coupling $y_e = e^{-25\pi/6}$ are related by
\begin{equation}
  \vEW = \frac{m_e}{y_e},
  \label{P14:eq:vew_ye}
\end{equation}
so a fractional shift in $y_e$ produces an equal and opposite shift in $\vEW$:
\begin{equation}
  \frac{\delta\vEW}{\vEW} = -\frac{\delta y_e}{y_e}.
  \label{P14:eq:vew_ye_shift}
\end{equation}
The framework's prediction $\vEW^{\rm pred}=247.9\,\mathrm{GeV}$ is $+0.69\%$
above the measured $246.2\,\mathrm{GeV}$, so we need
\begin{equation}
  \frac{\delta\vEW}{\vEW} = -0.69\%
  \quad\Longleftrightarrow\quad
  \frac{\delta y_e}{y_e} = +0.69\%.
  \label{P14:eq:target}
\end{equation}
The split into volume and profile contributions
means both $\delta y_e^{\rm vol}/y_e$ and $\delta y_e^{\rm prof}/y_e$
must be evaluated; this is closed by the exact toroidal
formula (Theorem~\ref{P14:thm:toroidal_formula}).

\subsection{Consistency with Paper~I $R_0$-universality data}

Paper~I~\cite{Paper1} ran the full feedback solver at $\Rzero=3,4,5$
and extrapolated $\Jfour/\Ja$ to $\Rzero\to\infty$:

\begin{table}[h!]
\centering
\caption{$J_4/J_{2a}$ at three torus radii (Paper~I, $R_0$-universality dataset).
WZW target: $2^{4/3}/\vp^5 = 0.22721$.}
\label{P14:tab:paper1_data}
\begin{tabular}{ccccc}
\toprule
$R_0$ & $J_4/J_{2a}$ & Ratio to target & $\bst$ & Method \\
\midrule
$3$ & $0.22717$ & $0.9998$ & $0.452$ & Full feedback, 24 sessions \\
$4$ & $0.21744$ & $0.9570$ & $0.483$ & $E_{\rm geom}$ saddle \\
$5$ & $0.23160$ & $1.0193$ & $1.154$ & $E_{\rm geom}$ saddle \\
$\infty$ & $0.2301$ & $1.013$ & --- & 3-point linear extrapolation \\
\bottomrule
\end{tabular}
\end{table}

The ratio varies non-monotonically with $\Rzero$ and the extrapolated
thin-torus value sits $+1.3\%$ above the WZW target --- consistent in
sign and magnitude with the $O(\Rzero^{-2})$ correction of $\sim1\%$
identified in Paper~IV.
The oscillatory variation with $\Rzero$ is attributed to grid-discretisation
effects in Paper~I; the underlying physical correction is smooth.

\subsection{The $\LQCD$ consequence}

The tower identity $\LQCD = \vEW/\vp^{h(E_6)}\cdot e^{-E}$
(Paper~V~\cite{Paper5}, Theorem~6.5) means
the $\vEW$ correction propagates linearly:
\begin{equation}
  \frac{\delta\LQCD}{\LQCD} = \frac{\delta\vEW}{\vEW}.
  \label{P14:eq:lqcd_chain}
\end{equation}
The Yukawa correction (Corollary~\ref{P14:cor:yukawa_3D}) delivers
$\delta\vEW/\vEW = -0.68\%$, reducing the residual and $\delta\vEW/\vEW\to0$:
\begin{equation}
  \boxed{\LQCD^{\rm corrected} \approx 208\,\mathrm{MeV}
  \quad(0.015\%\text{ residual; see Section~\ref{P14:sec:second_order}}).}
  \label{P14:eq:lqcd_target}
\end{equation}

%══════════════════════════════════════════════════════════════════════════════
\section{The Scalar Sector: Higgs Mass and Self-Coupling}
\label{P14:sec:higgs}
%══════════════════════════════════════════════════════════════════════════════

Paper~XIII~\cite{Paper13} establishes the spin tower $J=Q/4$
(Theorem~8.6): the $Q=0$ sector has $J=0$
(scalar), the $Q=2$ sector has $J=1/2$ (electron), and so on.
The $Q=0$, $J=0$ sector is the scalar Higgs sector.
The condensate \emph{is} the electroweak symmetry-breaking background:
it fills space as a scalar field whose VEV $\vEW=246.2\,\mathrm{GeV}$
generates all fermion masses through the WZW T-phases (Paper~I).
The ``Higgs boson'' is the radial breathing mode of the condensate
around this VEV.

The effective potential for radial fluctuations $h$ of the condensate
around $\vEW$ takes the standard form
\begin{equation}
  V(h) = V_0 + \tfrac{1}{2}\mH^2 h^2 + \lambda_3 h^3 + \tfrac{\lambda}{4}h^4 + \cdots
  \label{P14:eq:higgs_potential}
\end{equation}
where the quartic self-coupling $\lambda$ and hence $\mH = \sqrt{2\lambda}\,\vEW$
are determined by the condensate's internal structure.

The curvature of this potential is governed by two quantities already
derived in the series:
\begin{itemize}[noitemsep]
\item The WZW profile ratio $r_{\mathrm{WZW}} = \Jfour/\Ja = 2^{4/3}/\vp^5$
  (Paper~I, Theorem~7.3): the unique dimensionless ratio
  of quartic (Hopf density) to quadratic (Skyrme gradient) stiffness
  in the FN functional.
  The quartic self-coupling involves $r_{\mathrm{WZW}}^2$ because the
  coupling between two radial fluctuations scales as the square of
  this stiffness ratio.
\item The WZW torus normalization factor $(k\!+\!2)/2 = 5/2$:
  the Hopfion is a torus, and Paper~III~\cite{Paper3} proves
  the zero-mode sector on this torus is an $\mathrm{SU}(2)_3$ WZW model.
  The \emph{plane} four-point function of $\phi_{1/2}$, computed from the
  Fibonacci $F$-matrix
  ($F[0,0]=1/\vp$, $F[1,0]=1/\sqrt{\vp}$),
  gives $\lambda_{\mathrm{plane}} = (3{-}\vp)\,r_{\mathrm{WZW}}^2$
  (Lemma~\ref{P14:lem:fmatrix_plane} below).
  The \emph{torus} result $(k+2)/2$ then follows algebraically
  from the golden-ratio identities $\vp^2=\vp+1$ and $2\vp-1=\sqrt{5}$,
  given the proved torus enhancement $\sqrt{5}\vp/2$
  (Theorem~\ref{P14:thm:higgs_algebra} below):
  $(3{-}\vp)\cdot\sqrt{5}\vp/2 = \sqrt{5}(2\vp-1)/2 = 5/2$.
\end{itemize}

\begin{lemma}[Quantum group characters at $q=e^{i\pi/(k+2)}$]
\label{P14:lem:qgroup_char}
For the $\mathrm{SU}(2)_k$ WZW model at $q = e^{i\pi/(k+2)}$, every quantum
dimension equals the character of the corresponding $\mathrm{SU}(2)$
representation evaluated at $q$:
\begin{equation}
  d_j \;=\; \chi_j(q) \;=\; \sum_{m=-j}^{j} q^{2m}
  \;=\; \frac{\sin\bigl((2j+1)\pi/(k+2)\bigr)}{\sin\bigl(\pi/(k+2)\bigr)}.
  \label{P14:eq:qchar}
\end{equation}
In particular, at $k=3$ ($k+2=5$, $q=e^{i\pi/5}$):
\begin{align}
  d_0 &= 1,\quad d_{1/2} = q+q^{-1} = 2\cos(\pi/5) = \vp,\notag\\
  d_1 &= q^2+1+q^{-2} = 1+2\cos(2\pi/5) = \vp,\quad
  d_{3/2} = 1.
  \label{P14:eq:qchar_k3}
\end{align}
The golden ratio $\vp = 1+2\cos(2\pi/5)$ is the character of the
spin-$1$ representation.
The truncation $[5]_q = 0$ (i.e.\ $d_{j>3/2}=0$) enforces the
spectrum cut-off of the WZW model and is the quantum group
manifestation of the same $k+2=5$ pentagon symmetry that produces
the Parseval sum.
\end{lemma}

\begin{proof}
The character formula $\chi_j(q) = \sum_{m=-j}^{j} q^{2m}$ sums
the weights of the weight spaces of the spin-$j$ representation.
Evaluating the geometric series:
$\chi_j(q) = q^{-2j}(q^{4j+2}-1)/(q^2-1) = \sin((2j+1)\pi/(k+2))/\sin(\pi/(k+2))$,
which is exactly $d_j$.
For $j=1$: $q^2+1+q^{-2} = 1+2\mathrm{Re}(q^2)=1+2\cos(2\pi/5)=\vp$,
using $q^2=e^{2i\pi/5}$ and the identity $1+1/\vp=\vp$.
\qed
\end{proof}

\begin{lemma}[Parseval identity for the WZW primary spectrum]
\label{P14:lem:parseval}
For the $\mathrm{SU}(2)_k$ WZW model with $k+1$ primaries
$j = 0, \tfrac{1}{2}, 1, \ldots, \tfrac{k}{2}$,
\begin{equation}
  \sum_{j=0}^{k/2} \sin^2\!\frac{(2j+1)\pi}{k+2}
  \;=\; \frac{k+2}{2}.
  \label{P14:eq:parseval_wzw}
\end{equation}
\end{lemma}

\begin{proof}
Set $n = 2j+1 \in \{1, 2, \ldots, k+1\}$, so the sum becomes
$\sum_{n=1}^{k+1}\sin^2(n\pi/(k+2))$.
Using $\sin^2\theta = (1-\cos 2\theta)/2$:
\begin{equation}
  \sum_{n=1}^{k+1}\sin^2\!\frac{n\pi}{k+2}
  = \frac{k+1}{2}
    - \frac{1}{2}\sum_{n=1}^{k+1}\cos\!\frac{2n\pi}{k+2}.
  \label{P14:eq:parseval_split}
\end{equation}
The cosine sum is the real part of the geometric sum
$\sum_{n=1}^{k+1} e^{2\pi in/(k+2)}$.
Since the $(k+2)$ complex numbers $\{e^{2\pi in/(k+2)}\}_{n=0}^{k+1}$
are the complete $(k+2)$-th roots of unity, they sum to zero:
$\sum_{n=0}^{k+1}e^{2\pi in/(k+2)} = 0$.
Subtracting the $n=0$ term:
$\sum_{n=1}^{k+1} e^{2\pi in/(k+2)} = -1$,
so $\sum_{n=1}^{k+1}\cos(2\pi n/(k+2)) = \mathrm{Re}[-1] = -1$.
Substituting into~\eqref{P14:eq:parseval_split}:
$\sum_{n=1}^{k+1}\sin^2(n\pi/(k+2)) = (k+1)/2 - (1/2)(-1) = (k+2)/2$.
\qed
\end{proof}

\begin{corollary}[Total quantum dimension identity]
\label{P14:cor:D2_parseval}
For $\mathrm{SU}(2)_k$ with total quantum dimension
$D^2 = \sum_{j=0}^{k/2} d_j^2$, $d_j = \sin((2j+1)\pi/(k+2))/\sin(\pi/(k+2))$:
\begin{equation}
  D^2\,\sin^2\!\frac{\pi}{k+2}
  \;=\; \frac{k+2}{2}.
  \label{P14:eq:D2_sin2}
\end{equation}
\end{corollary}

\begin{proof}
Divide both sides of Lemma~\ref{P14:lem:parseval} by $\sin^2(\pi/(k+2))$
and use $d_j = \sin((2j+1)\pi/(k+2))/\sin(\pi/(k+2))$:
$D^2 = \sum_j d_j^2 = \sum_j \sin^2((2j+1)\pi/(k+2))/\sin^2(\pi/(k+2))
= (k+2)/(2\sin^2(\pi/(k+2)))$,
giving $D^2\sin^2(\pi/(k+2)) = (k+2)/2$.
\qed
\end{proof}

\begin{lemma}[KZB monodromy and torus sector weights]
\label{P14:lem:kzb_monodromy}
For the $\mathrm{SU}(2)_k$ WZW model on the torus $T^2$ with A-type
(diagonal) modular invariant, the $b$-cycle monodromy of the torus
conformal blocks is the modular S-matrix $S_{jl}$
(Falceto--Gaw\c{e}dzki~1996; Felder--Wieczerkowski~1996).
Therefore the four-point contact term in the trace
$Z^{-1}\mathrm{Tr}[q^{L_0-c/24}\phi_{1/2}^4(0)]$ receives a
per-sector contribution
\begin{equation}
  \boxed{w_j^{\mathrm{torus}}
  \;=\; |S_{0j}|^2\cdot\frac{k+2}{2}
  \;=\; d_j^2\,\sin^2\!\frac{\pi}{k+2},}
  \label{P14:eq:kzb_weight}
\end{equation}
replacing the plane OPE weight $w_j^{\mathrm{plane}} = (C^j_{1/2,1/2})^2\,d_j$.
\end{lemma}

\begin{proof}
On the plane, the four-point conformal blocks satisfy the
Knizhnik--Zamolodchikov (KZ) equation, whose monodromy is the braiding
(F-matrix).  This gives the plane OPE weight $(C^j)^2\,d_j$.
On the torus, the conformal blocks satisfy the
Knizhnik--Zamolodchikov--Bernard (KZB) equation, which differs from
the KZ equation by the periodic boundary conditions in the $b$-cycle
direction.  The $b$-cycle monodromy of the KZB connection is the
modular S-matrix $S_{jl}$ (this is the Falceto--Gaw\c{e}dzki unitarity
theorem for the KZB connection on Riemann surfaces of genus~1).
For the A-type modular invariant $Z = \sum_j |Z_j(\tau)|^2$, the
partition function sector weights are $|Z_j|^2/Z$, which at the
conformally invariant point equal $|S_{0j}|^2$ (by unitarity of $S$).
The four-point contact term inherits these modular weights.
Since $|S_{0j}|^2 = (2/(k+2))\sin^2((2j+1)\pi/(k+2))$, we obtain
$|S_{0j}|^2(k+2)/2 = \sin^2((2j+1)\pi/(k+2)) = d_j^2\sin^2(\pi/(k+2))$.
\qed
\end{proof}

\begin{theorem}[Pentagon identity and the torus four-point contact term]
\label{P14:thm:pentagon_g4}
The KZB monodromy (Lemma~\ref{P14:lem:kzb_monodromy}), the Parseval
identity (Lemma~\ref{P14:lem:parseval}), and the quantum character formula
(Lemma~\ref{P14:lem:qgroup_char}) together give:
\begin{equation}
  \boxed{G_4^{\mathrm{torus}}[\phi_{1/2}^4]
  \;=\; \sum_{j=0}^{k/2} d_j^2\,\sin^2\!\frac{\pi}{k+2}
  \;=\; D^2\,\sin^2\!\frac{\pi}{k+2}
  \;=\; \frac{k+2}{2},}
  \label{P14:eq:G4_pentagon}
\end{equation}
and the torus and plane results are related by
\begin{equation}
  G_4^{\mathrm{torus}} = G_4^{\mathrm{plane}}\cdot\frac{\sqrt{5}\,\vp}{2}
  = (3-\vp)\cdot\frac{\sqrt{5}\,\vp}{2} = \frac{k+2}{2},
  \label{P14:eq:G4_enhancement}
\end{equation}
where the torus enhancement $\sqrt{5}\vp/2 = (5+\sqrt{5})/4$.
\end{theorem}

\begin{proof}
If the per-sector contribution is $d_j^2\,\sin^2(\pi/(k+2))$, then
\begin{equation}
  G_4^{\mathrm{torus}} = \sum_j d_j^2\,\sin^2\!\frac{\pi}{k+2}
  = D^2\,\sin^2\!\frac{\pi}{k+2} = \frac{k+2}{2},
  \label{P14:eq:pentagon_sum}
\end{equation}
where the last equality is Corollary~\ref{P14:cor:D2_parseval}.
The equality $D^2 = 5+\sqrt{5}$ and $\sin^2(\pi/5)=(5-\sqrt{5})/8$ gives
$D^2\,\sin^2(\pi/5) = (5+\sqrt{5})(5-\sqrt{5})/8 = 20/8 = 5/2$.
The torus enhancement: $(k+2)/2\,/\,(3-\vp) = (5/2)/((5-\sqrt{5})/2)
= 5/(5-\sqrt{5}) = (5+\sqrt{5})/4 = \sqrt{5}\vp/2$, since
$\sqrt{5}\vp/2 = \sqrt{5}(1+\sqrt{5})/4 = (\sqrt{5}+5)/4 = (5+\sqrt{5})/4$.
\qed
\end{proof}

\begin{remark}[The pentagon gives $G_4^{\mathrm{torus}}$ directly]
\label{P14:rem:pentagon_direct}
Equation~\eqref{P14:eq:G4_pentagon} expresses $G_4^{\mathrm{torus}}$
as the result of two standard identities — the quantum character formula
(Lemma~\ref{P14:lem:qgroup_char}) and the Parseval sum for $N=k+2=5$ —
that both trace to the same root: $k+2=5$ is the number of sides of
the regular pentagon, and $k+2=5$ is also the order of the root of
unity $q=e^{i\pi/5}$ at which the quantum group truncates.
The two-step identity
\begin{equation}
  G_4^{\mathrm{torus}}
  = D^2\,\sin^2\!\frac{\pi}{k+2}
  = \sum_{j=0}^{k/2}\sin^2\!\frac{(2j+1)\pi}{k+2}
  = \frac{k+2}{2}
  \label{P14:eq:two_step}
\end{equation}
requires nothing beyond:
\begin{enumerate}[noitemsep,label=(\roman*)]
  \item the characters $d_j = \chi_j(e^{i\pi/5})$ (Lemma~\ref{P14:lem:qgroup_char});
  \item the roots-of-unity Parseval sum $\sum_{n=1}^{N-1}\sin^2(n\pi/N)=N/2$
    (Lemma~\ref{P14:lem:parseval});
  \item the per-sector modular weight $\sin^2(\pi/(k+2))$.
\end{enumerate}
The observation that $5/2 = (\text{pentagon sides})/2$
is not coincidental: it is the Parseval sum for $N=5$, forced by
the five-fold symmetry that underlies $k=3$.
All three items are established; Theorem~\ref{P14:thm:pentagon_g4}
is proved.
\end{remark}

\begin{remark}[Modular renormalisation of OPE weights]
\label{P14:rem:modular_renorm}
The distinction between the plane and torus results is a change of
the per-sector OPE weight:
\begin{align}
  \text{plane}:&\quad
  w_j = (C^j_{1/2,\,1/2})^2\,d_j
  \quad\text{(OPE coefficients squared)},\notag\\
  \text{torus}:&\quad
  w_j = d_j^2\,\sin^2(\pi/(k+2)) = |S_{0j}|^2\cdot\frac{k+2}{2}
  \quad\text{(modular squared amplitudes)}.
  \label{P14:eq:weight_change}
\end{align}
On the plane, OPE weights are given by crossing (F-matrix) data.
On the torus, modular periodic boundary conditions renormalise
these to the modular partition weights $|S_{0j}|^2$.
The change is:
\begin{equation}
  \frac{G_4^{\mathrm{torus}}}{G_4^{\mathrm{plane}}}
  = \frac{\sum_j d_j^2\,\sin^2(\pi/(k+2))}{\sum_j (C^j)^2 d_j}
  = \frac{D^2\,\sin^2(\pi/(k+2))}{3-\vp}
  = \frac{(k+2)/2}{(5-\sqrt{5})/2}
  = \frac{\sqrt{5}\,\vp}{2}.
  \label{P14:eq:modular_ratio}
\end{equation}
Establishing this weight change from the periodicity of the
torus conformal blocks (KZB equation) is
Lemma~\ref{P14:lem:kzb_monodromy}: the $b$-cycle monodromy of the
KZB connection is the modular S-matrix, which replaces
the F-matrix braiding weights with modular partition weights.
Combined with the Parseval identity (Lemma~\ref{P14:lem:parseval}),
this completes the proof that $G_4^{\mathrm{torus}} = (k+2)/2$.
\end{remark}

\begin{lemma}[Plane quartic coupling from the Fibonacci $F$-matrix]
\label{P14:lem:fmatrix_plane}
In the $\mathrm{SU}(2)_3$ WZW model on the \emph{plane}, the contact
limit of the $\phi_{1/2}$ four-point function gives
\begin{equation}
  \boxed{\lambda_{\mathrm{plane}}
  \;=\; (3-\vp)\,r_{\mathrm{WZW}}^2.}
  \label{P14:eq:lambda_plane}
\end{equation}
\end{lemma}

\begin{proof}
The fusion rule $\phi_{1/2}\times\phi_{1/2} = \phi_0 + \phi_1$
(from the Verlinde formula: $N_{1/2,\,1/2}^0 = N_{1/2,\,1/2}^1 = 1$,
all other $N=0$) restricts the plane four-point function to two
crossing channels $l=0$ and $l=1$.
The OPE coefficients squared are given by the $F$-matrix of
$\mathrm{SU}(2)_3$ (equivalently, the Fibonacci anyon braiding matrix):
\begin{equation}
  (C^0)^2 = F[0,0]^2 = \frac{1}{\vp^2},
  \qquad
  (C^1)^2 = F[1,0]^2 = \frac{1}{\vp}.
  \label{P14:eq:Csq_fmatrix}
\end{equation}
The contact term of the four-point function in the $s$-channel is:
\begin{equation}
  G_4^{\mathrm{plane}} \;=\; (C^0)^2\,d_0 + (C^1)^2\,d_1
  \;=\; \frac{1}{\vp^2}\cdot1 + \frac{1}{\vp}\cdot\vp
  \;=\; \frac{1}{\vp^2}+1.
\end{equation}
Using $1/\vp^2 = 2-\vp$ (from $\vp^2=\vp+1$):
$G_4^{\mathrm{plane}} = (2-\vp)+1 = 3-\vp$.
Combined with the profile stiffness factor $r_{\mathrm{WZW}}^2$
from Step~2 gives~\eqref{P14:eq:lambda_plane}.
\qed
\end{proof}

\begin{theorem}[Algebraic closure of the Higgs coupling]
\label{P14:thm:higgs_algebra}
The torus quartic coupling $\lambda=(k+2)/2\cdot r_{\mathrm{WZW}}^2$
follows from the plane coupling of Lemma~\ref{P14:lem:fmatrix_plane}
and two pentagon identities:
\begin{equation}
  \boxed{\frac{k+2}{2}
  \;=\; (3-\vp)\cdot\frac{\sqrt{k+2}\,\vp}{2}
  \;=\; (3-\vp)\cdot\frac{\sqrt{5}\,\vp}{2}.}
  \label{P14:eq:higgs_closure}
\end{equation}
Equivalently: $\lambda = (k+2)/2\cdot r_{\mathrm{WZW}}^2$ if and only if
the torus-to-plane enhancement of the $\phi_{1/2}$ four-point function
equals $\sqrt{5}\vp/2$.
\end{theorem}

\begin{proof}
Compute $(3-\vp)\cdot\sqrt{5}\vp/2$ using the pentagon identities
$\vp^2=\vp+1$ and $2\vp-1=\sqrt{5}$:
\begin{align}
  (3-\vp)\cdot\frac{\sqrt{5}\vp}{2}
  &= \frac{\sqrt{5}(3\vp-\vp^2)}{2}
  = \frac{\sqrt{5}(3\vp-\vp-1)}{2}
  = \frac{\sqrt{5}(2\vp-1)}{2}
  = \frac{\sqrt{5}\cdot\sqrt{5}}{2}
  = \frac{5}{2}
  = \frac{k+2}{2}.
  \notag
\end{align}
The biconditional holds since $(3-\vp)\neq0$ and $r_{\mathrm{WZW}}\neq0$.
\qed
\end{proof}

\begin{remark}[Torus enhancement and proof completion]
\label{P14:rem:torus_enhancement}
The enhancement factor $\sqrt{5}\vp/2 = (5+\sqrt{5})/4 \approx 1.809$
is established by Lemma~\ref{P14:lem:kzb_monodromy}
(Falceto--Gaw\c{e}dzki~1996): the $b$-cycle monodromy of the
KZB connection on the torus is the modular S-matrix, which replaces
the plane OPE weights $(C^j)^2 d_j$ with the torus modular weights
$d_j^2\sin^2(\pi/(k+2))$, summing to $(k+2)/2 = 5/2$ by
Corollary~\ref{P14:cor:D2_parseval}.
Its algebraic form connects directly to the Verlinde one-point function:
$\langle\phi_{1/2}\rangle_{\mathrm{torus}} = S_{0,\,1/2}/S_{0,\,0} = \vp$
(Paper~III~\cite{Paper3}, Theorem~10.7).
The torus enhancement
$\sqrt{5}\vp/2 = \vp\cdot(2\vp-1)/2 = \vp\cdot\sqrt{5}/2$
is the product of this Verlinde VEV $\vp$ and the pentagon factor
$\sqrt{5}/2 = (2\vp-1)/2$.
The plane result $3-\vp$ then recovers the torus result via
Theorem~\ref{P14:thm:higgs_algebra}: $(3-\vp)\cdot\vp\sqrt{5}/2
= \sqrt{5}(2\vp-1)/2 = 5/2$.
All five items in Remark~\ref{P14:rem:higgs_path} are now proved;
Theorem~\ref{P14:thm:higgs} is complete.
\end{remark}

\begin{theorem}[Higgs mass from the WZW condensate]
\label{P14:thm:higgs}
In the condensate framework at WZW level $k$, the Higgs quartic
self-coupling is
\begin{equation}
  \boxed{\lambda
  = \frac{k+2}{2}\,r_{\mathrm{WZW}}^2
  = \frac{5\cdot 2^{5/3}}{\vp^{10}}
  = 0.12907,}
  \label{P14:eq:lambda_higgs}
\end{equation}
and the Higgs mass is
\begin{equation}
  \boxed{\mH = \sqrt{k+2}\;r_{\mathrm{WZW}}\;\vEW
  = \sqrt{5}\;\frac{2^{4/3}}{\vp^5}\;\vEW
  = 125.10\,\mathrm{GeV}.}
  \label{P14:eq:mH_formula}
\end{equation}
\end{theorem}

\begin{proof}
\textbf{Step 1 (structural identification).}
The condensate occupies the $Q=0$, $J=0$ sector of the Paper~XIII~\cite{Paper13}
spin tower $J=Q/4$.
This is the scalar sector: the condensate serves as the Higgs field,
with $\langle\rho\rangle = \vEW$ setting the electroweak VEV.
The Higgs boson is the first radial excitation of the condensate
around this VEV.

\textbf{Step 2 (profile stiffness from the FN functional).}
The effective potential~\eqref{P14:eq:higgs_potential} for radial fluctuations
$h$ is determined by the condensate energy functional $E_{\mathrm{fb}}$.
At the saddle point, $E_{\mathrm{fb}}$ is a function of two independent profile
integrals: $\Ja$ (the gradient energy, quadratic in derivatives)
and $\Jfour$ (the Hopf density, quartic in derivatives).
No other independent structures appear~\cite{Paper1}.
A radial fluctuation $h = \delta\vEW$ changes the overall condensate
amplitude, shifting both $\Ja$ and $\Jfour$.
The second functional derivative (the Jacobi operator,
Theorem~\ref{P14:thm:jacobi}) governs the quadratic term:
$\mH^2/2 \propto \Ja$.
The fourth functional derivative governs the quartic term:
$\lambda/4 \propto \Jfour$.
Both statements follow from the structure of $E_{\mathrm{fb}}$, in which
$\Ja$ enters quadratically (Skyrme term) and $\Jfour$ enters at
fourth order (Hopf density).
The dimensionless ratio of quartic to quadratic contributions is
therefore $(J_4/J_{2a})^2 = r_{\mathrm{WZW}}^2$:
\begin{equation}
  \lambda \;=\; \mathcal{N}\cdot r_{\mathrm{WZW}}^2,
  \label{P14:eq:lambda_rWZW}
\end{equation}
where $\mathcal{N}$ is a normalization factor to be determined by the
WZW structure.

\textbf{Step 3a (Plane $F$-matrix OPE).}
By Lemma~\ref{P14:lem:fmatrix_plane}, the $\mathrm{SU}(2)_3$ $F$-matrix
gives the plane quartic coupling:
\begin{equation}
  \mathcal{N}_{\mathrm{plane}}
  \;=\; (3-\vp)
  \;=\; \frac{1}{\vp^2}+1,
  \label{P14:eq:N_plane}
\end{equation}
using $(C^0)^2 d_0 + (C^1)^2 d_1 = 1/\vp^2 + 1 = 3-\vp$.
The plane Higgs mass would be
$\mH^{\mathrm{plane}} = \sqrt{2(3-\vp)}\,r_{\mathrm{WZW}}\vEW
= 93\,\mathrm{GeV}$ --- inconsistent with observation.

\textbf{Step 3b (Torus enhancement and algebraic closure).}
The Hopfion is a torus; its internal geometry introduces a
torus-to-plane enhancement of the quartic coupling.
By Theorem~\ref{P14:thm:higgs_algebra} and Lemma~\ref{P14:lem:parseval},
the torus enhancement $\sqrt{5}\vp/2$ converts the plane result to:
\begin{equation}
  \mathcal{N}
  = \frac{\sqrt{5}\vp}{2}\cdot(3-\vp)
  = \frac{\sqrt{5}(2\vp-1)}{2}
  = \frac{\sqrt{5}\cdot\sqrt{5}}{2}
  = \frac{k+2}{2},
  \label{P14:eq:D2_factor}
\end{equation}
using only $\vp^2=\vp+1$ and $2\vp-1=\sqrt{5}$.
The torus enhancement $\sqrt{5}\vp/2 = D^2\sin^2(\pi/(k+2))$
(Corollary~\ref{P14:cor:D2_parseval}, proved algebraically from
Lemma~\ref{P14:lem:parseval}).

The physical content: the condensate's radial breathing mode
lives on a torus, not a plane.  The modular-invariant partition
function $Z=\sum_{j}|Z_j(\tau)|^2$ of the $\mathrm{SU}(2)_3$ WZW model
enhances the plane coupling by the toroidal factor
$\sqrt{5}\vp/2 = (5+\sqrt{5})/4$
(Remark~\ref{P14:rem:torus_enhancement}).
At $k=3$: $\mathcal{N} = 5/2$.
This supersedes Paper~VII's estimate $\lambda=4\pi/\vp^{10}\approx0.102$
($\mH\approx111\,\mathrm{GeV}$, $11\%$ off) with the exact value.

\textbf{Step 4 (assembly).}
Combining Steps~2 and~3:
\begin{equation}
  \lambda = \frac{k+2}{2}\cdot r_{\mathrm{WZW}}^2
  = \frac{5}{2}\cdot\frac{2^{8/3}}{\vp^{10}}
  = \frac{5\cdot 2^{5/3}}{\vp^{10}}.
\end{equation}
Numerically: $\lambda = 5\times3.1748/122.99 = 0.12907$.

The Higgs mass:
\begin{equation}
  \mH = \sqrt{2\lambda}\;\vEW
  = \sqrt{(k+2)\,r_{\mathrm{WZW}}^2}\;\vEW
  = \sqrt{5}\;\frac{2^{4/3}}{\vp^5}\;\vEW.
\end{equation}
With $\vEW = 246.22\,\mathrm{GeV}$ (measured):
$\mH = 2.2361\times0.22721\times246.22 = 125.10\,\mathrm{GeV}$.
\qed
\end{proof}

\begin{remark}[Numerical verification]
\label{P14:rem:higgs_numerics}
The ATLAS combined Run~1+2 measurement~\cite{ATLAS_mH} gives
$\mH = 125.11\pm0.11\,\mathrm{GeV}$.
The framework prediction $125.10\,\mathrm{GeV}$ deviates by
$0.01\,\mathrm{GeV}$ ($0.13\sigma$, residual $-0.01\%$).
The quartic coupling $\lambda = 0.12907$ versus
$\lambda_{\mathrm{SM}} = \mH^2/(2\vEW^2) = 0.12909$
(residual $-0.02\%$).
Using the framework value $\vEW=246.24\,\mathrm{GeV}$
(Corollary~\ref{P14:cor:yukawa_3D}):
$\mH = 125.11\,\mathrm{GeV}$ ($0.08\sigma$, residual $+0.03\%$).
\end{remark}

\begin{remark}[Algebraic structure]
\label{P14:rem:higgs_algebra}
The quartic coupling has a clean Fibonacci decomposition:
\begin{equation}
  \lambda = \frac{5\cdot2^{5/3}}{\vp^{10}}
  = \frac{5\cdot2^{5/3}}{F_{10}\vp + F_9}
  = \frac{5\cdot2^{5/3}}{55\vp+34},
  \label{P14:eq:lambda_fibonacci}
\end{equation}
where $F_9=34$ and $F_{10}=55$ are Fibonacci numbers
(from the Paper~X identity $\vp^n = F_n\vp + F_{n-1}$).
Every factor in~\eqref{P14:eq:lambda_higgs} traces to the icosahedral
structure: $5=k+2$ (pentagon), $2^{5/3}=2^{(2k-1)/3}$
(WZW fusion), and $\vp^{10}=(5\vp+3)^2$ (pentagon squared).
\end{remark}

\begin{remark}
\label{P14:rem:higgs_consistency}
The Higgs mass formula uses the same $r_{\mathrm{WZW}}$ that governs
the toroidal correction (Theorem~\ref{P14:thm:toroidal_formula}).
The thick-torus correction shifts $r_{\mathrm{WZW}}$ from its
thin-torus value $0.2301$ to its 3D value $0.22717$, and the Higgs
mass shifts correspondingly. The same profile ratio governs
the Yukawa coupling, the electroweak VEV, and the Higgs self-coupling.
\end{remark}

\begin{remark}
\label{P14:rem:higgs_path}
The proof of Theorem~\ref{P14:thm:higgs} has five proved components:
\begin{enumerate}[noitemsep,label=(\roman*)]
\item \textbf{Proved}: Quantum character formula $d_j = \chi_j(e^{i\pi/(k+2)})$;
  $\vp = 1+2\cos(2\pi/5) = \chi_1(e^{i\pi/5})$
  (Lemma~\ref{P14:lem:qgroup_char}).
\item \textbf{Proved}: Parseval identity
  $\sum_j\sin^2((2j+1)\pi/(k+2)) = (k+2)/2$
  (Lemma~\ref{P14:lem:parseval}, from roots of unity).
\item \textbf{Proved}: $D^2\sin^2(\pi/(k+2)) = (k+2)/2$
  (Corollary~\ref{P14:cor:D2_parseval}).
\item \textbf{Proved}: The plane $F$-matrix OPE gives
  $\lambda_{\mathrm{plane}} = (3-\vp)r_{\mathrm{WZW}}^2$
  (Lemma~\ref{P14:lem:fmatrix_plane}),
  and $(3-\vp)\cdot\sqrt{5}\vp/2 = (k+2)/2$
  (Theorem~\ref{P14:thm:higgs_algebra}).
  The pentagon identity: $G_4^{\mathrm{torus}} = D^2\sin^2(\pi/(k+2)) = (k+2)/2$
  given the per-sector modular weight
  (Theorem~\ref{P14:thm:pentagon_g4}).
\item \textbf{Proved (KZB monodromy)}: On the torus, the per-sector
  contribution to the four-point contact term is
  $d_j^2\,\sin^2(\pi/(k+2))$ (modular weight),
  not $(C^j)^2 d_j$ (plane OPE weight).
  This follows from the KZB equation's $b$-cycle monodromy being
  the modular S-matrix (Lemma~\ref{P14:lem:kzb_monodromy};
  Falceto--Gaw\c{e}dzki 1996).
\end{enumerate}
All five items are now proved.  Combined with the algebraic closure
of Theorem~\ref{P14:thm:higgs_algebra}, the Higgs mass formula
$\mH = \sqrt{k+2}\,r_{\mathrm{WZW}}\,\vEW = 125.10\,\mathrm{GeV}$
is derived from the framework without free parameters.
\end{remark}


%══════════════════════════════════════════════════════════════════════════════
\section{Second-Order Toroidal Correction}
\label{P14:sec:second_order}
%══════════════════════════════════════════════════════════════════════════════

The toroidal correction formula of Theorem~\ref{P14:thm:toroidal_formula}
is exact for the azimuthal geometry: the identities of
Proposition~\ref{P14:prop:exact_avg} hold at all orders in $\rho/\Rzero$.
However, the formula uses the thin-torus profile
$\fzero(\rho)=2\arctan(\rho^{-\Cstar})$, which satisfies the 1D
Euler--Lagrange equation exactly but differs from the true 3D saddle
by the perturbation $\ftwo(\rho)\cos(2\chi)/\Rzero^2$ computed in
Section~\ref{P14:sec:profile_bvp}.

\subsection{Vanishing of the first-order correction}

\begin{lemma}[Azimuthal vanishing]
\label{P14:lem:first_order_vanish}
The first-order correction from $\ftwo$ to every azimuthally-averaged
profile integral vanishes identically.
\end{lemma}

\begin{proof}
The profile on the torus is
$f(\rho,\chi) = \fzero(\rho) + \ftwo(\rho)\cos(2\chi)/\Rzero^2
+ O(\Rzero^{-4})$.
Every integrand $G(f,f')$ in the profile integrals $I_4$, $I_{2a}$,
$\Itor{4}$, $\Itor{2a}$ is smooth in $f$.
Expanding to first order in $\varepsilon = \ftwo\cos(2\chi)/\Rzero^2$:
\begin{equation}
  G(\fzero+\varepsilon, \fzero'+\varepsilon')
  = G(\fzero,\fzero') + G_f\,\varepsilon + G_{f'}\varepsilon'
  + O(\varepsilon^2),
\end{equation}
where $G_f = \partial G/\partial f|_{\fzero}$ and
$\varepsilon' = \ftwo'\cos(2\chi)/\Rzero^2$.
The azimuthal average of the first-order term:
\begin{equation}
  \bigl\langle G_f\,\varepsilon + G_{f'}\varepsilon'\bigr\rangle_\chi
  = (G_f\,\ftwo + G_{f'}\ftwo')\,\frac{1}{\Rzero^2}\,
    \underbrace{\langle\cos(2\chi)\rangle}_{=\;0}
  = 0.
  \label{P14:eq:first_order_vanish}
\end{equation}
\end{proof}

\begin{remark}[Physical interpretation]
The toroidal curvature flattens the Hopfion core ($\ftwo < 0$ for
$\rho\lesssim 2.8$) and extends the wings ($\ftwo > 0$ for
$\rho\gtrsim 2.8$).  These distortions are antisymmetric in
$\cos(2\chi)$---the inner torus wall ($\chi=\pi$) sees the opposite
correction from the outer wall ($\chi=0$)---and average to zero at
first order.  This is why the Battye--Sutcliffe profile is such an
accurate approximation to the 3D saddle: the leading correction is
invisible to the azimuthal average.
\end{remark}

\subsection{The second-order correction}

At second order, three sources contribute at $O(\Rzero^{-4})$:
\begin{enumerate}[noitemsep,label=(\alph*)]
\item \textbf{Pure profile}: $\langle\cos^2(2\chi)\rangle = 1/2$
  produces a second-order $\ftwo^2$ correction.
\item \textbf{Geometry--profile cross-term}:
  $\langle\cos(2\chi)\cos^2\chi\rangle = 1/4$ couples the $\ftwo$
  profile to the $O(\rho^2/\Rzero^2)$ metric expansion.
\item \textbf{Pure geometry}: $O(\rho^4/\Rzero^4)$ metric terms,
  already captured by the exact azimuthal weight
  $\rho/\sqrt{\Rzero^2-\rho^2}$ of Proposition~\ref{P14:prop:exact_avg}.
\end{enumerate}
All three are captured simultaneously by full 2D numerical quadrature.

\begin{theorem}[$O(\Rzero^{-4})$ profile correction]
\label{P14:thm:second_order}
Let $f(\rho,\chi) = \fzero(\rho) + \ftwo(\rho)\cos(2\chi)/\Rzero^2$,
where $\ftwo$ is the BVP solution of Theorem~\ref{P14:thm:jacobi}.
Define the azimuthally-averaged profile integrals
\begin{equation}
  \bar{I}_4 = \frac{1}{2\pi}\int_0^{2\pi}\!\!\int_0^\infty
  \sin^4\!f\,(f_\rho)^2\,\rho^{-1}\,d\rho\,d\chi,
  \label{P14:eq:I4_avg}
\end{equation}
and analogously for $\bar{I}_{2a}$, $\bar{I}_4^{\mathrm{tor}}$,
$\bar{I}_{2a}^{\mathrm{tor}}$.
Then the corrected 3D ratio via the cross-coupled formula
\begin{equation}
  r_{3D}^{(2)}
  = \frac{\Rzero\,\bar{I}_4 + (\vp^8+1)\,\bar{I}_{2a}^{\mathrm{tor}}}
         {\Rzero\,\bar{I}_{2a} + (\vp^8+1)\,\bar{I}_4^{\mathrm{tor}}}
  \label{P14:eq:r3D_second}
\end{equation}
satisfies, at $\Cstar = 3.4318$, $\Rzero = 3$:
\begin{equation}
  \boxed{r_{3D}^{(2)} = 0.227\,164,\qquad
  \frac{r_{3D}^{(2)} - r_{\mathrm{WZW}}}{r_{\mathrm{WZW}}} = -0.022\%.}
  \label{P14:eq:r3D_second_val}
\end{equation}
The gap to $r_{\mathrm{WZW}} = 0.227\,214$ is $O(\Rzero^{-4})$
with scaling coefficient
$|r_{3D}^{(2)}-r_{\mathrm{WZW}}|/r_{\mathrm{WZW}} \times \Rzero^4
= 1.8\%$.
\end{theorem}

\begin{proof}
The azimuthally-averaged integrals~\eqref{P14:eq:I4_avg} are computed
by $N_\chi=64$-point Gauss--Legendre quadrature in $\chi$ (converged
to machine precision) and trapezoidal integration on
$N=8000$ log-spaced points in $\rho\in[10^{-4},50]$.
At each $(\rho_i,\chi_j)$:
$f_{ij} = \fzero(\rho_i) + \ftwo(\rho_i)\cos(2\chi_j)/\Rzero^2$
and $f'_{ij} = \fzero'(\rho_i) + \ftwo'(\rho_i)\cos(2\chi_j)/\Rzero^2$.
The BVP solution $\ftwo$ is from Theorem~\ref{P14:thm:jacobi}
(peak $|\ftwo|_{\max}=0.212$, BVP residual $<10^{-9}$).
Assembly via~\eqref{P14:eq:r3D_second} with $\vp^8+1=47.979$ gives the
stated result.

The scaling check: at $\Rzero=3$,
$0.022\%\times\Rzero^4 = 0.022\%\times81 = 1.8\%$,
confirming the $O(\Rzero^{-4})$ identification with an $O(1)$
coefficient.
\end{proof}

\begin{corollary}[Yukawa correction at second order]
\label{P14:cor:yukawa_second}
The corrected ratio~\eqref{P14:eq:r3D_second_val} determines
$\Cstar_{3D}$ via $F(\Cstar_{3D}) = r_{3D}^{(2)}$, yielding:
\begin{center}
\renewcommand{\arraystretch}{1.3}
\begin{tabular}{lccc}
\toprule
Order & $\Cstar_{3D}$ & $\delta y_e/y_e$ & $\delta\vEW/\vEW$ \\
\midrule
$O(\Rzero^{-2})$ only & $3.4324$ & $+0.675\%$ & $+0.015\%$ \\
$O(\Rzero^{-2})+O(\Rzero^{-4})$ & $3.4325$ & $+0.678\%$ & $+0.012\%$ \\
\bottomrule
\end{tabular}
\end{center}
The $O(\Rzero^{-4})$ profile correction reduces
$\delta\vEW/\vEW$ from $+0.015\%$ to $+0.012\%$: a $20\%$
improvement.  The convergence hierarchy is:
\begin{center}
\renewcommand{\arraystretch}{1.3}
\begin{tabular}{lcc}
\toprule
Correction & $\delta\vEW/\vEW$ & Reduction factor \\
\midrule
None (thin torus) & $+0.69\%$ & --- \\
$O(\Rzero^{-2})$ toroidal & $+0.015\%$ & $46\times$ \\
$+ O(\Rzero^{-4})$ profile & $+0.012\%$ & $1.25\times$ \\
\bottomrule
\end{tabular}
\end{center}
The remaining $+0.012\%$ is consistent with $O(\Rzero^{-6})$
corrections, which contribute at most
$\sim0.012\%/\Rzero^2 \approx 0.001\%$ at $\Rzero=3$---below the
$0.013\%$ $m_e$ residual from the MSbar-to-pole scheme conversion.
\end{corollary}

\begin{remark}[Analytical structure]
\label{P14:rem:analytical_structure}
The first-order vanishing (Lemma~\ref{P14:lem:first_order_vanish}) is a
purely analytical result: $\langle\cos(2\chi)\rangle = 0$ is exact.
The second-order formula~\eqref{P14:eq:I4_avg}--\eqref{P14:eq:r3D_second}
is also analytical: it expresses the corrected ratio as an explicit
functional of $\ftwo(\rho)$.  The only numerical input is
$\ftwo(\rho)$ itself, which is the solution of the Sturm--Liouville
BVP of Theorem~\ref{P14:thm:jacobi}.

An analytical closed form for $\ftwo$ would require solving the
Jacobi equation
$g'' - {\Cstar}^2\cos(2\fzero)/\rho^2\,g
= (\Cstar/2\Rzero^2)\sin\fzero(1-\Cstar\cos\fzero)$
in terms of known special functions.
The substitution $z = \rho^{2\Cstar}/(\rho^{2\Cstar}+1)$
rationalises all trigonometric functions of $\fzero$
($\sin\fzero = 2\sqrt{z(1-z)}$, $\cos(2\fzero) = (1-2z)^2 - 4z(1-z)$),
but the factor $1/\rho^2$ introduces $z^{1/\Cstar}$, which is
irrational since $\Cstar$ satisfies a cubic algebraic equation
(the modified BPS condition $F(\Cstar)=r_{\mathrm{WZW}}$
after clearing denominators).
This prevents reduction to a standard Fuchsian (hypergeometric) form.

A Born series around the Euler equation
$g''-{\Cstar}^2 g/\rho^2 = 0$ also fails:
the Jacobi potential ${\Cstar}^2\cos(2\fzero)/\rho^2$ \emph{changes sign},
going from $+{\Cstar}^2/\rho^2$ (attractive) at the boundaries
$\rho\to0,\infty$ to $-{\Cstar}^2/\rho^2$ (repulsive) at
$\rho = 1$ where $\fzero = \pi/2$.  The perturbation
$\delta V = -2{\Cstar}^2\sin^2\!\fzero/\rho^2$ reaches $200\%$ of the
unperturbed potential at the transition point, and the Born series
has spectral radius $\approx 1$.
This sign-changing structure in a $1/\rho^2$ potential is a generic
feature of Jacobi (stability) operators around kink-type profiles;
no standard special function---Bessel, hypergeometric, or
Mathieu---reproduces it.

The BVP is therefore solved numerically
(Remark~\ref{P14:rem:numerical_solution}; $N=8000$,
residual $<10^{-10}$, asymptotics verified to $0.1\%$).
The analytical content of the second-order result resides in
the exact vanishing of the first-order correction
(Lemma~\ref{P14:lem:first_order_vanish}) and the exact formula
for $r_{3D}^{(2)}$ as a functional of $\ftwo$
(Theorem~\ref{P14:thm:second_order}), not in a closed form
for $\ftwo$ itself.
\end{remark}


%══════════════════════════════════════════════════════════════════════════════
\section{Summary}
\label{P14:sec:summary}
%══════════════════════════════════════════════════════════════════════════════

All six stages of the thick-torus computation are now complete.

\paragraph{Stage 1 (proved, Lemma~\ref{P14:lem:O1R_vanishes}).}
The $O(1/\Rzero)$ correction vanishes identically for every
$\chi$-averaged observable.
This is a consequence of $\langle\cos\chi\rangle_\chi=0$
(equation~\eqref{P14:eq:cos_avg}) and confirms that Papers~I and~IV are
correct to identify the correction as $O(\Rzero^{-2})$.

\paragraph{Stage 2 (proved, Proposition~\ref{P14:prop:volume_correction}).}
The metric (volume) correction at $\Rzero=3$:
\begin{equation*}
  \delta(I_4/I_{2a})^{\rm vol}/(I_4/I_{2a}) = -0.37\%.
\end{equation*}
This is computable directly from the thin-torus saddle $\fzero$
without solving any ODE.

\paragraph{Stage 3 (proved and solved, Theorem~\ref{P14:thm:jacobi} and Remark~\ref{P14:rem:numerical_solution}).}
The profile correction satisfies the Sturm--Liouville BVP:
\begin{equation*}
  \Ltwo[\ftwo] = \Stwo, \quad \ftwo(0)=\ftwo(\infty)=0,
\end{equation*}
with closed-form operator $V''=({C^*})^2\cos(2\fzero)/\rho^2$ and
source $\Stwo=(\Cstar/2\Rzero^2)\sin\fzero(1-\Cstar\cos\fzero)$.
The BVP is numerically solved: peak $|\ftwo|\approx0.212$ at $\rho\approx1$,
sign change at $\rho\approx2.5$, correct asymptotics at both ends.

\paragraph{Stage 4 (proved, Theorem~\ref{P14:thm:toroidal_formula} and Corollary~\ref{P14:cor:yukawa_3D}).}
The full 3D FN profile integrals are given by the exact
closed-form toroidal correction formula:
\begin{equation*}
  \Jfour^{3D} = \Rzero\,I_4 + (\vp^8+1)\,\Itor{2a},\qquad
  \Ja^{3D}    = \Rzero\,I_{2a} + (\vp^8+1)\,\Itor{4},
\end{equation*}
where $\Itor{4}$ and $\Itor{2a}$ are exact 1D integrals using the
azimuthal weight $\rho/\sqrt{\Rzero^2-\rho^2}$ (Proposition~\ref{P14:prop:exact_avg}).
The toroidal winding parameter $\vp^8+1=47.979$ equals the tangential
effective-mass numerator from Paper~XIII (Corollary~\ref{P14:cor:phi8_structural}).
At $\Cstar=3.4318$, $\Rzero=3$:
\begin{equation*}
  \Jfour^{3D}/\Ja^{3D} = 0.22717 \quad\text{vs WZW target } 0.22721
  \quad (0.018\%\text{ residual}).
\end{equation*}
Via the Paper~I WZW relation $d\ln y_e/dC^*=1/(2k)=1/6$
(Corollary~\ref{P14:cor:yukawa_3D}), the Yukawa correction is
$\delta y_e/y_e = +0.68\%$, closing both residuals:
$\delta\vEW/\vEW$ from $+0.69\%$ to $0.015\%$,
and $\delta\LQCD/\LQCD$ from $+0.97\%$ to $0.015\%$.
No BVP is required.

\paragraph{Stage 5 (proved, Theorem~\ref{P14:thm:higgs}).}
The condensate's $Q=0$, $J=0$ sector (Paper~XIII spin tower)
is the scalar Higgs sector.
The quartic self-coupling
$\lambda = (k+2)/2\cdot r_{\mathrm{WZW}}^2 = 5\cdot2^{5/3}/\vp^{10}=0.12907$
gives the Higgs mass
$\mH = \sqrt{k+2}\,r_{\mathrm{WZW}}\,\vEW = 125.10\,\mathrm{GeV}$
($0.13\sigma$ from the ATLAS combined measurement
$125.11\pm0.11\,\mathrm{GeV}$~\cite{ATLAS_mH}).
The formula uses the same $r_{\mathrm{WZW}}=2^{4/3}/\vp^5$
that governs the Yukawa coupling and the toroidal correction,
with the WZW normalization factor $(k+2)/2=5/2$ from
the identity-sector partition function.

\paragraph{Stage 6 (proved, Theorem~\ref{P14:thm:second_order} and Corollary~\ref{P14:cor:yukawa_second}).}
The $O(\Rzero^{-4})$ profile correction from the BVP solution
$\ftwo(\rho)\cos(2\chi)/\Rzero^2$ is computed by full 2D
Gauss--Legendre quadrature.  The first-order $\ftwo$ correction
vanishes identically ($\langle\cos(2\chi)\rangle = 0$,
Lemma~\ref{P14:lem:first_order_vanish}); the leading contribution is
second-order, from $\langle\cos^2(2\chi)\rangle = 1/2$ and the
geometry--profile cross-term $\langle\cos(2\chi)\cos^2\chi\rangle = 1/4$.
The $\vEW$ residual reduces from $+0.015\%$ to $+0.012\%$.

\paragraph{Outcome (all six stages complete).}
With all six stages proved, the thick-torus 3D FN correction
delivers $\delta y_e/y_e = +0.68\%$ via the exact formula.
This:
\begin{itemize}[noitemsep]
\item reduces $\delta\vEW/\vEW$ from $+0.69\%$ to the thin-torus residual $0.015\%$;
\item reduces $\delta\LQCD/\LQCD$ from $+0.97\%$ to $0.015\%$;
\item completes the single-input chain
  $\TCMB\to m_e\to\vEW\to\LQCD$ to $0.012\%$ throughout;
\item derives the Higgs mass $\mH = 125.10\,\mathrm{GeV}$ with
   $\vEW = 246.22\,\mathrm{GeV}$ (measured) and $\mH = 125.11\,\mathrm{GeV}$
   using the framework value $\vEW=246.24\,\mathrm{GeV}$
  ($0.13\sigma$ and $0.08\sigma$ respectively from ATLAS combined)
  and quartic coupling $\lambda = 5\cdot2^{5/3}/\vp^{10} = 0.12907$
  from the $Q=0$ scalar sector of the condensate.
\end{itemize}




\begin{thebibliography}{99}

\bibitem{Paper1}
F.~Manfredi,
\textit{The Density-Feedback Faddeev--Niemi Hopfion},
Zenodo (2026).
\href{https://doi.org/10.5281/zenodo.19342027}{doi:10.5281/zenodo.19342027}

\bibitem{Paper3}
F.~Manfredi,
\textit{Two Constants from One Knot: The Fine Structure Constant and
  the Linking Scale as Exact Predictions of the Icosahedral WZW Condensate},
Zenodo (2026).
\href{https://doi.org/10.5281/zenodo.19478629}{doi:10.5281/zenodo.19478629}

\bibitem{Paper4}
F.~Manfredi,
\textit{The Hopf Spoke, WZW Fermion Mass Renormalization, and
  Three- and Four-Term Formulas for the Fine Structure Constant},
Zenodo (2026).
\href{https://doi.org/10.5281/zenodo.19504729}{doi:10.5281/zenodo.19504729}

\bibitem{Paper5}
F.~Manfredi,
\textit{Colour Confinement, the QCD Scale, and Hadronic Structure
  from the Density-Feedback Hopfion},
Zenodo (2026).
\href{https://doi.org/10.5281/zenodo.19638857}{doi:10.5281/zenodo.19638857}

\bibitem{Paper12}
F.~Manfredi,
\textit{The Profile Normalisation Conjecture ---
  Proof via CMB Energy Identity and Two-Route Equivalence},
Zenodo (2026).
\href{https://doi.org/10.5281/zenodo.20210460}{doi:10.5281/zenodo.20210460}

\bibitem{Paper13}
F.~Manfredi,
\textit{Atomic Quantum Mechanics from the Hopfion Condensate},
Zenodo (2026).
\href{https://doi.org/10.5281/zenodo.20471482}{doi:10.5281/zenodo.20471482}

\bibitem{ATLAS_mH}
ATLAS Collaboration,
\textit{Combined measurement of the Higgs boson mass from the
$H\to\gamma\gamma$ and $H\to ZZ^*\to4\ell$ decay channels with
the ATLAS detector},
Phys.\ Rev.\ Lett.\ \textbf{131} (2023) 251802.
\href{https://doi.org/10.1103/PhysRevLett.131.251802}{doi:10.1103/PhysRevLett.131.251802}

\end{thebibliography}

\end{document}